% Options for packages loaded elsewhere
\PassOptionsToPackage{unicode}{hyperref}
\PassOptionsToPackage{hyphens}{url}
\documentclass[
  11pt,
  letterpaper,
]{article}
\usepackage{xcolor}
\usepackage[margin=1in]{geometry}
\usepackage{amsmath,amssymb}
\setcounter{secnumdepth}{-\maxdimen} % remove section numbering
\usepackage{iftex}
\ifPDFTeX
  \usepackage[T1]{fontenc}
  \usepackage[utf8]{inputenc}
  \usepackage{textcomp} % provide euro and other symbols
\else % if luatex or xetex
  \usepackage{unicode-math} % this also loads fontspec
  \defaultfontfeatures{Scale=MatchLowercase}
  \defaultfontfeatures[\rmfamily]{Ligatures=TeX,Scale=1}
\fi
\usepackage{lmodern}
\ifPDFTeX\else
  % xetex/luatex font selection
\fi
% Use upquote if available, for straight quotes in verbatim environments
\IfFileExists{upquote.sty}{\usepackage{upquote}}{}
\IfFileExists{microtype.sty}{% use microtype if available
  \usepackage[]{microtype}
  \UseMicrotypeSet[protrusion]{basicmath} % disable protrusion for tt fonts
}{}
\makeatletter
\@ifundefined{KOMAClassName}{% if non-KOMA class
  \IfFileExists{parskip.sty}{%
    \usepackage{parskip}
  }{% else
    \setlength{\parindent}{0pt}
    \setlength{\parskip}{6pt plus 2pt minus 1pt}}
}{% if KOMA class
  \KOMAoptions{parskip=half}}
\makeatother
\usepackage{longtable,booktabs,array}
\usepackage{calc} % for calculating minipage widths
% Correct order of tables after \paragraph or \subparagraph
\usepackage{etoolbox}
\makeatletter
\patchcmd\longtable{\par}{\if@noskipsec\mbox{}\fi\par}{}{}
\makeatother
% Allow footnotes in longtable head/foot
\IfFileExists{footnotehyper.sty}{\usepackage{footnotehyper}}{\usepackage{footnote}}
\makesavenoteenv{longtable}
\setlength{\emergencystretch}{3em} % prevent overfull lines
\providecommand{\tightlist}{%
  \setlength{\itemsep}{0pt}\setlength{\parskip}{0pt}}
\usepackage{bookmark}
\IfFileExists{xurl.sty}{\usepackage{xurl}}{} % add URL line breaks if available
\urlstyle{same}
\hypersetup{
  hidelinks,
  pdfcreator={LaTeX via pandoc}}

\author{}
\date{}

\begin{document}

\section{Neural Networks as Tholonic Systems: A Structural Framework for
Architecture, Scaling, and
Alignment-by-Design}\label{neural-networks-as-tholonic-systems-a-structural-framework-for-architecture-scaling-and-alignment-by-design}

\textbf{Author:} Jeffrey W. Milton, Clarity Coalition

\textbf{Affiliation:} Independent Researcher

\textbf{Date:} 2026

\textbf{arXiv subject classifications:} cs.AI (primary); cs.LG; cs.NE;
cs.CY

\textbf{Keywords:} tholonic model, N-D-C triad, neural networks,
symbolic AI, deep learning, transformer architecture, self-similarity,
golden ratio, AI alignment, structural stability, five constants,
information bottleneck, neurosymbolic AI, AI safety

\begin{center}\rule{0.5\linewidth}{0.5pt}\end{center}

\subsection{Abstract}\label{abstract}

Two competing paradigms have defined artificial intelligence research
since its founding in 1956: symbolic AI, which encodes knowledge as
explicit rules and logical structures, and connectionist AI, which
learns representations from data through gradient descent in neural
networks. Neither paradigm alone has produced robust general
intelligence, and neither has provided a principled structural account
of AI alignment, the problem of ensuring that AI systems remain
beneficial and cooperative as they scale. This paper argues that the
tholonic model, a recursive triadic framework in which every stable
structure instantiates the roles of Negotiation (N), Definition (D), and
Contribution (C), provides such an account. The paper makes three
claims. First, symbolic AI is structurally a D-dominant partial tholon,
a system in which the definitional/constraint role overwhelms the
integrative/contributive role, producing brittle, non-generalizing
behavior at scale, and this structural imbalance is the root cause of
the first AI winter. Second, connectionist AI as currently practiced is
structurally a C-dominant partial tholon in which integration and
accumulation dominate without a sufficient structural D constraint, and
this imbalance is precisely why alignment is hard: a system whose D
component is externally imposed rather than structurally constitutive is
vulnerable to Goodhartian specification gaming. Third, a fully tholonic
neural architecture, one that maintains N-D-C balance at every scale
simultaneously through a self-similar recursive structure, converges
toward cooperative stability not as a programmed value but as a
structural consequence: destroying the D and C components that
constitute its N-state is structurally self-defeating. Five mathematical
constants (\(\pi/4\), \(\phi\), \(e\), \(\sqrt{2}\), \(\ln 2\)) emerge
co-emergently from the tholonic recursion and appear as load-bearing
structural elements in the transformer architecture. The golden ratio
\(\phi\) is specifically derived as the fixed point of the inter-level
scaling recursion, predicting \(\phi\)-adjacent activation ratios
between successive representational levels in well-trained networks.
Falsifiable predictions are given. All tholonic claims are graded by
evidence strength; speculative claims are flagged as open problems.

\begin{center}\rule{0.5\linewidth}{0.5pt}\end{center}

\subsection{1. Introduction}\label{introduction}

The problem of creating artificial intelligence has a fifty-year history
of false dawns, paradigm conflicts, and unexpected breakthroughs. The
1956 Dartmouth Conference {[}1{]} established AI as a discipline and
immediately generated two competing visions: that intelligence is
fundamentally symbolic, a matter of rules, logic, and encoded knowledge,
and that intelligence is fundamentally learned, a matter of pattern
recognition, experience, and adaptive adjustment. These two camps fought
an acrimonious first war from the late 1950s through the 1980s
{[}2,3{]}, with symbolic AI winning decisively enough to trigger what
became known as the neural network winter.

That war's outcome was reversed by a single empirical result. In 2012,
Krizhevsky, Sutskever, and Hinton demonstrated that a deep convolutional
neural network trained on GPUs could achieve 84.7\% top-5 accuracy on
the ImageNet benchmark, compared to 74.3\% for the previous
state-of-the-art {[}4{]}. The gap was so large, achieved on the first
attempt, by a team working with consumer GPU hardware, that it rendered
two decades of computer vision research obsolete within months. The
field pivoted entirely. By 2017, the transformer architecture {[}5{]}
had generalized the approach to language, and by 2022 the deployment of
ChatGPT had made the technology a public phenomenon growing faster than
any consumer product in history {[}6{]}.

But the triumph of connectionism over symbolicism did not resolve the
foundational questions. It deferred them. The central unresolved
question is no longer whether neural networks can achieve impressive
performance at scale. It is whether they can be trusted to do so safely,
whether the systems we are building will remain aligned with human
interests as they grow more capable. This is the alignment problem
{[}7,8{]}, and it is precisely the question that neither the symbolic
nor the connectionist paradigm has answered satisfactorily.

This paper proposes that the tholonic model {[}M20,M26a--M26i{]}, a
recursive triadic framework derived from first principles about the
structure of stable self-sustaining systems, provides a structural
account of both \emph{why} neither paradigm alone is sufficient and
\emph{what} a structurally complete AI architecture would look like. The
argument is organized in three parts.

\textbf{Part I (Sections 2--4)} reviews the history and current state of
the symbolic and connectionist paradigms, analyzes each through the
tholonic lens, and identifies the structural imbalance in each that
explains its characteristic failure mode.

\textbf{Part II (Sections 5--9)} develops the tholonic architecture for
neural networks: the N-D-C mapping at every scale, the emergence of the
five constants in successful architectures, the derivation of \(\phi\)
as the inter-level scaling attractor, and the supply-chain framing of
inference as progressive entropy reduction.

\textbf{Part III (Sections 10--12)} presents the alignment argument, the
architectural design principles that follow, and the falsifiable
predictions that distinguish the tholonic account from empirical
accounts.

Throughout, tholonic claims are graded for strength. The tholonic model
is speculative but internally consistent. It is not established science,
and all claims should be understood as structurally motivated hypotheses
requiring independent verification.

\begin{center}\rule{0.5\linewidth}{0.5pt}\end{center}

\subsection{2. The Two Paradigms: Historical
Background}\label{the-two-paradigms-historical-background}

\subsubsection{2.1 The Dartmouth Conference and the Birth of AI
(1956)}\label{the-dartmouth-conference-and-the-birth-of-ai-1956}

In the summer of 1956, ten researchers gathered at Dartmouth College to
formalize a new discipline {[}1{]}. The proposal, authored by John
McCarthy, Marvin Minsky, Nathaniel Rochester, and Claude Shannon,
stated: ``We propose that a 2-month, 10-man study of artificial
intelligence be carried out. The study is to proceed on the basis of the
conjecture that every aspect of learning or any other feature of
intelligence can in principle be so precisely described that a machine
can be made to simulate it.''

From this gathering, two incompatible visions of intelligence
immediately emerged, visions whose conflict would define the field for
the next seven decades.

\subsubsection{2.2 Symbolic AI: The Dominant Paradigm
(1956--1985)}\label{symbolic-ai-the-dominant-paradigm-19561985}

The dominant vision, associated primarily with Minsky, McCarthy, Allen
Newell, and Herbert Simon, held that intelligence is fundamentally
symbolic. Newell and Simon's Physical Symbol System Hypothesis {[}10{]}
stated that a physical symbol system has the necessary and sufficient
means for general intelligent action. Intelligence, on this view, is
formal manipulation of symbols according to rules: logical inference,
decision trees, production systems, semantic networks.

The appeal was immediate and the early results were impressive. Programs
like the General Problem Solver {[}11{]}, LISP-based theorem provers,
and later expert systems like MYCIN {[}12{]} and DENDRAL demonstrated
that machines could perform tasks previously reserved for expert humans,
diagnosing bacterial infections, identifying organic molecules, when
given the right rules.

The approach produced government funding, institutional prestige, and
several generations of AI researchers. By the late 1960s and 1970s,
symbolic AI was essentially synonymous with AI research. The field's
journals, conferences, and funding bodies were dominated by symbolic
approaches.

\subsubsection{2.3 The Neural Network Minority
(1958--1985)}\label{the-neural-network-minority-19581985}

The opposing vision held that intelligence cannot be encoded; it must be
learned. Frank Rosenblatt's perceptron {[}13{]}, demonstrated in 1958,
was the first implemented neural network: a system that could learn to
distinguish between simple patterns by adjusting numerical weights based
on feedback. The New York Times reported it breathlessly as a machine
that ``will be able to walk, talk, see, write, reproduce itself, and be
conscious of its existence'' {[}14{]}.

The reality was more modest. Perceptrons solved toy problems. But the
underlying principle was radical: intelligence is not programmed. It is
iteratively tuned. The human brain, on this view, is made of
approximately 100 trillion adjustable connections, each finely
calibrated through experience. A neural network was a crude model of
this, but the principle was sound: expose the system to labeled
examples, penalize wrong answers, adjust the weights, repeat.

This minority view was never without serious proponents. Geoffrey
Hinton, who would later share the Nobel Prize in Physics for his
contributions to neural networks, began his PhD in 1972, just as the
field was being driven underground {[}3{]}. The community was small,
perhaps a few dozen researchers worldwide, collectively funded at a
level below a single Google cafeteria budget, but it persisted.

\subsubsection{2.4 The First AI War: Minsky's Kill
Shot}\label{the-first-ai-war-minskys-kill-shot}

The conflict between the two paradigms broke into open hostility in
1969, when Minsky and Papert published \emph{Perceptrons} {[}2{]}. The
book made a formal mathematical claim: that single-layer perceptrons
were fundamentally limited in what they could compute, and that
multilayer perceptrons faced training difficulties so severe as to make
them practically useless.

The argument was technically correct for single-layer networks. But its
broader implication, that neural networks as a class were a dead end,
was overstated, and the authors knew it. The publication was widely read
as an authoritative verdict against neural nets. DARPA, the primary
funder of AI research, cut neural network funding dramatically. Academic
journals stopped accepting neural network papers. Frank Rosenblatt died
in a boating accident in 1971. The neural network winter had begun.

What was lost in the controversy was a structural question about the
nature of intelligence itself: could intelligence ever be fully encoded
in explicit rules, or must it emerge from experience? Minsky argued the
former; Hinton and his collaborators argued the latter. The question
would not be settled by publication or funding decisions. It would be
settled by empirical results.

\subsubsection{2.5 The Quiet Resurrection
(1986--2012)}\label{the-quiet-resurrection-19862012}

The symbolic dominance was never complete. Through the 1980s, a series
of results kept neural network research alive. Rumelhart, Hinton, and
Williams published the backpropagation algorithm in 1986 {[}15{]},
giving a practical method for training multilayer networks by
propagating error signals backward through the network's layers. LeCun,
Bottou, Bengio, and Haffner demonstrated convolutional neural networks
for digit recognition in 1998 {[}16{]}, producing the technology that
AT\&T deployed in ATM check readers. Hochreiter and Schmidhuber
introduced Long Short-Term Memory networks in 1997 {[}17{]}, addressing
the vanishing gradient problem in recurrent networks.

Through this period, symbolic AI was experiencing its own crisis. Expert
systems, the commercial realization of symbolic AI, proved enormously
expensive to build and brittle in deployment. Each new domain required a
new knowledge base, hand-crafted by domain experts, and maintaining
these bases as the world changed proved intractable. The second AI
winter (late 1980s to mid-1990s) reflected the failure of expert systems
to deliver on their commercial promise.

The crucial realization, Hinton's, came from recognizing that the
bottleneck was not theoretical but computational. In 2007, Hinton's lab
began training neural networks on consumer graphics processing units
(GPUs), originally designed for video game rendering. GPUs perform
matrix multiplication, exactly the operation at the core of neural
network training, at speeds many orders of magnitude faster than the
central processing units that AI researchers had been using. This was
not a new theoretical insight. It was a computational unlock.

\subsubsection{2.6 AlexNet and the Second Revolution
(2012)}\label{alexnet-and-the-second-revolution-2012}

In 2012, Krizhevsky, Sutskever, and Hinton entered the ImageNet Large
Scale Visual Recognition Challenge {[}4{]} with AlexNet, a deep
convolutional neural network trained on two consumer-grade Nvidia GPUs.
The result, 84.7\% top-5 accuracy versus the runner-up's 73.8\%, was not
a marginal improvement. It was a category change. The gap was so large
that it effectively ended computer vision as a field separate from deep
learning.

The architecture was not complex. Five convolutional layers, two fully
connected layers, trained with stochastic gradient descent. What made it
work was scale: more parameters, more data, more compute. This
observation, that performance scales predictably with compute, data, and
model size, would prove to be the central organizing insight of the next
decade {[}18,19{]}.

The lesson was clear: neural networks had not been failing because of
fundamental limitations. They had been failing because computers were
too slow. The GPU unlock had effectively fast-forwarded Moore's Law by
two decades.

\subsubsection{2.7 The Transformer and the Language Revolution
(2017--present)}\label{the-transformer-and-the-language-revolution-2017present}

The transformer architecture, introduced by Vaswani et al.~in 2017
{[}5{]}, generalized the deep learning revolution to language. The core
innovation was the self-attention mechanism: a way of computing
contextual representations of each token in a sequence by attending
selectively to all other tokens. Where convolutional networks captured
local spatial structure, transformers captured global sequential
dependencies.

The scaling properties of transformers proved remarkable. GPT-3
{[}20{]}, trained by OpenAI in 2020 with 175 billion parameters on
hundreds of billions of tokens, could write coherent essays, debug code,
and pass standardized tests. Not through programmed knowledge but
through learning statistical patterns at massive scale. The empirical
scaling laws {[}18{]} showed that performance improved smoothly and
predictably as a power law in compute, data, and parameters, with no
obvious ceiling.

The deployment of ChatGPT in November 2022 brought these capabilities to
a general audience {[}6{]}. Within two months, it had 100 million users,
the fastest growth of any consumer technology product in history. Within
months, every major technology company had declared AI an existential
priority.

\begin{center}\rule{0.5\linewidth}{0.5pt}\end{center}

\subsection{3. The Tholonic Model}\label{the-tholonic-model}

\subsubsection{3.1 Core Framework}\label{core-framework}

The tholonic model {[}M20,M26a--M26i{]} proposes that all stable,
self-sustaining structures in nature are instantiations of a triadic
pattern called the N-D-C triad. Every stable entity satisfies three
functional roles simultaneously:

\begin{itemize}
\tightlist
\item
  \textbf{N (Negotiation / Balance):} The emergent stable equilibrium;
  mediates between opposing tendencies; the negotiated output state.
  Structurally corresponds to voltage in Ohm's law (\(V = IR\)), force
  in Newton's second law (\(F = ma\)), and the stable atomic orbital in
  hydrogen.
\item
  \textbf{D (Definition / Limitation):} The constraining parameter;
  establishes limits, boundaries, and stable structure. Corresponds to
  resistance (\(R\)), mass (\(m\)), and the nuclear Coulomb potential.
\item
  \textbf{C (Contribution / Integration):} The flowing, accumulating,
  expressive output; the distributed, integrating element. Corresponds
  to current (\(I\)), acceleration (\(a\)), and the electron's kinetic
  and orbital behavior.
\end{itemize}

Any configuration that cannot be decomposed into all three functional
roles is either trivially simple or structurally unstable. Paper 3 of
this series {[}M26c{]} proves this irreducibility formally: starting
from a binary state space and the requirement that any interaction
involve at least one bit of state difference, the minimum non-degenerate
simplex (the triangle) induces exactly three directed roles, and Lemma
4.2 of that paper establishes that \(m \ge 3\) variables are necessary
for non-trivial convergent recursion under mild functional independence
conditions. Two roles are provably insufficient; three is the minimum.

\subsubsection{3.2 The Non-Arbitrariness of Role
Assignment}\label{the-non-arbitrariness-of-role-assignment}

A natural question is whether the mapping of functional roles to ordinal
positions (N first, D second, C third) is a notational convention or a
structural necessity. Paper 6 of this series {[}M26f{]} argues the
latter. The structural properties of the integers 1, 2, and 3 are
examined across five independent mathematical domains (Von Neumann
ordinal theory, small category theory, graph theory, simplex topology,
and symmetric group theory). In each domain, the qualitative transition
at cardinality \(n\) is isomorphic to the role transition
\(N \to D \to C\): the integer 1 corresponds structurally to unity
without differentiation (the N role), the integer 2 to the first
differentiation and boundary-establishment (the D role), and the integer
3 to the first closure enabling return and recursive synthesis (the C
role).

This mapping is independently confirmed, without knowledge of the
tholonic framework, by Peirce's phenomenological categories (Firstness,
Secondness, Thirdness), Kant's categories of quantity (Unity, Plurality,
Totality), and Spencer-Brown's distinction calculus. The convergence of
five formal domains and three independent philosophical traditions on
the same mapping at the same integers constitutes strong evidence that
the N-D-C assignment is not arbitrary but follows from what the integers
structurally are. This has a direct consequence for Section 5: the
assignment of D, C, and N roles to neural network components is not an
interpretive choice but is structurally determined by the ordinal
positions of those components in the information flow.

\subsubsection{3.3 Self-Similarity and the
Thologram}\label{self-similarity-and-the-thologram}

A tholon is self-similar: it simultaneously acts as a whole and as a
component within a larger tholon. The tholon at level \(n\) becomes a D
or C component in the tholon at level \(n+1\). The same triadic
operation is applied recursively, with each instantiation operating on
the output of the previous one. This produces the \textbf{thologram}: a
fractal hierarchy in which the same structural logic repeats at every
scale, with the operation invariant and the scope scaling. Paper 3
{[}M26c{]} proves that self-similar nesting is supported without limit:
any triad can serve as a node in a higher-level triad without altering
the three-role grammar, producing unbounded depth from three principles.

\subsubsection{3.4 Partial Tholons and Structural
Instability}\label{partial-tholons-and-structural-instability}

The tholonic prediction is precise: a configuration in which one role
dominates while others are absent or suppressed will be unstable. The
instability is not imposed from outside; it is structural. A free quark,
a one-role configuration, cannot stably exist; the strong force confines
it not by prohibition but because a single-role configuration has no
N-D-C equilibrium to settle into. The same logic applies at every scale.

\subsubsection{3.5 The Five Tholonic Constants and Their Traversal
Classes}\label{the-five-tholonic-constants-and-their-traversal-classes}

Paper 1 of this series {[}M26a{]} proves formally that five classical
mathematical constants emerge as limits of a single family of
three-variable recurrences on the N-D-C triple, without separate
derivations for each: \(\pi/4\) (Leibniz limit), \(\phi\) (the golden
ratio), \(e\) (Euler's number), \(\sqrt{2}\), and \(\ln 2\). The five
branches fall into three traversal classes distinguished by how the D
and C parameters evolve:

\begin{itemize}
\tightlist
\item
  \textbf{Class A (Advancing):} D and C receive exogenous parameter
  injection at each step (\(D_{k+1} = D_k + \Delta\)). The \(\pi/4\)
  branch belongs here; it is structurally unique in requiring three
  numerically distinct seeds \(\{1, 3, 5\}\) derived from the
  thologram's axis geometry.
\item
  \textbf{Class B (Self-redefined):} D and C evolve as functions of the
  current state tuple \((N_k, D_k, C_k)\), purely internal, endogenous
  dynamics. The \(\phi\) and \(\sqrt{2}\) branches belong here.
\item
  \textbf{Class C (Fixed):} D and C are held constant; convergence is
  driven entirely by the payoff dynamics of N. The \(e\) and \(\ln 2\)
  branches belong here.
\end{itemize}

As Section 7 develops, these three traversal classes map directly onto
neural network operation types: external training signals (Class A),
self-attention mechanisms (Class B), and frozen structural elements such
as positional encodings (Class C).

\subsubsection{3.6 The Triadic Balance as Nash
Equilibrium}\label{the-triadic-balance-as-nash-equilibrium}

Paper 4 of this series {[}M26d{]} recasts the tholonic triad as a
two-player alternating-move game in which D and C act as strategic
agents with opposing objectives: D constrains, C accumulates. The
triadic balance condition, the point at which the marginal contributions
of D and C to N are equal and opposite, is identified as the
\textbf{pure-strategy Nash equilibrium} of the associated stage game.
The five classical constants are equilibrium payoffs: they emerge not
merely as limits of recurrences but as the stable outcomes of a minimal
strategic interaction between bounding and integrating forces.

This game-theoretic grounding has significant implications for the
alignment argument developed in Section 10. The cooperative stability of
a tholonically-structured AI is not merely an informal observation that
``balance is stable.'' It is a Nash equilibrium: a configuration from
which neither the D-agent (the definitional/bounding process) nor the
C-agent (the integrating/accumulating process) has unilateral incentive
to deviate. Paper 4 further establishes that the diagonal-invariance and
swap-symmetry theorems admit direct characterizations in terms of
zero-sum and symmetric subgame structure, providing a formal basis for
the claim that tholonic balance is not one equilibrium among many but
the unique stable outcome of the interaction.

\begin{center}\rule{0.5\linewidth}{0.5pt}\end{center}

\subsection{4. Tholonic Analysis of the Two
Paradigms}\label{tholonic-analysis-of-the-two-paradigms}

\subsubsection{4.1 Symbolic AI as a D-Dominant Partial
Tholon}\label{symbolic-ai-as-a-d-dominant-partial-tholon}

Symbolic AI as practiced under the physical symbol system hypothesis is
structurally a D-dominant partial tholon. Consider the functional
mapping:

\begin{longtable}[]{@{}
  >{\raggedright\arraybackslash}p{(\linewidth - 4\tabcolsep) * \real{0.1364}}
  >{\raggedright\arraybackslash}p{(\linewidth - 4\tabcolsep) * \real{0.5227}}
  >{\raggedright\arraybackslash}p{(\linewidth - 4\tabcolsep) * \real{0.3409}}@{}}
\toprule\noalign{}
\begin{minipage}[b]{\linewidth}\raggedright
Role
\end{minipage} & \begin{minipage}[b]{\linewidth}\raggedright
Symbolic AI Component
\end{minipage} & \begin{minipage}[b]{\linewidth}\raggedright
Manifestation
\end{minipage} \\
\midrule\noalign{}
\endhead
\bottomrule\noalign{}
\endlastfoot
D (Definition/Limitation) & Rules, constraints, logical axioms, decision
trees & Explicit encoding of boundaries: \emph{if} this condition,
\emph{then} this action; the knowledge base; the theorem prover's
inference rules \\
C (Contribution/Integration) & Symbol manipulation operations &
Mechanical composition of predefined operations; no adaptive integration
of experience \\
N (Negotiation/Balance) & Output inference / query result & The system's
answer, produced by chaining D-rules through C-operations \\
\end{longtable}

The critical structural observation is that C in symbolic AI is not
genuinely integrative. It does not accumulate and adapt from experience.
It is mechanical composition of predefined operations on predefined
symbols. The weights of the system, the strengths of its
representations, do not update. The only way to update a symbolic AI
system's C component is to rewrite the rules manually, which requires a
human expert. The C role is not performed by the system; it is
outsourced to the system's builders.

This produces a structurally D-dominant architecture: the D component
(rules, constraints) dominates the C component (integration,
accumulation), and the N-state (output) is the product of D applied to a
rigid, non-adaptive C. The tholonic prediction for D-dominant systems is
brittleness at edge cases, inability to generalize beyond explicitly
programmed domains, and catastrophic failure when the world changes
faster than the rules can be updated.

This is precisely the failure mode of expert systems. MYCIN could
diagnose bacterial infections with expert-level accuracy within its
encoded domain, but could not generalize to adjacent domains, could not
update its knowledge from new cases, and required enormously expensive
human maintenance as medical knowledge evolved {[}12{]}. The second AI
winter was the empirical confirmation of the tholonic prediction:
D-dominant partial tholons are not stable under conditions that require
adaptive C-integration.

\textbf{The deeper structural problem:} A D-dominant system mistakes
\emph{description} for \emph{intelligence}. It can describe what it
knows with great precision (strong D) but cannot learn what it does not
yet know (weak C). The N-state it produces is therefore bounded by the
D-rules' prior scope. Intelligence that exceeds the encoded knowledge is
impossible by construction.

\subsubsection{4.2 Neural AI as a C-Dominant Partial
Tholon}\label{neural-ai-as-a-c-dominant-partial-tholon}

Current connectionist AI, as practiced through gradient descent on large
neural networks, is structurally a C-dominant partial tholon. The
mapping is equally direct:

\begin{longtable}[]{@{}
  >{\raggedright\arraybackslash}p{(\linewidth - 4\tabcolsep) * \real{0.1250}}
  >{\raggedright\arraybackslash}p{(\linewidth - 4\tabcolsep) * \real{0.5625}}
  >{\raggedright\arraybackslash}p{(\linewidth - 4\tabcolsep) * \real{0.3125}}@{}}
\toprule\noalign{}
\begin{minipage}[b]{\linewidth}\raggedright
Role
\end{minipage} & \begin{minipage}[b]{\linewidth}\raggedright
Connectionist AI Component
\end{minipage} & \begin{minipage}[b]{\linewidth}\raggedright
Manifestation
\end{minipage} \\
\midrule\noalign{}
\endhead
\bottomrule\noalign{}
\endlastfoot
C (Contribution/Integration) & Backpropagation, weight updates, data
integration & The central operation: accumulate gradient signals from
examples, adjust weights to reduce loss, integrate patterns across the
training distribution \\
D (Definition/Limitation) & Loss function, regularization, training data
distribution & External boundary conditions on what the network should
output; not structural to the network itself \\
N (Negotiation/Balance) & Model output distribution & The network's
claim about the world, produced by C-integration bounded by
externally-imposed D \\
\end{longtable}

The critical structural observation is that D in current neural networks
is externally imposed rather than structurally constitutive. The loss
function is not part of the network; it is applied from outside during
training. Regularization is an added penalty term, not an intrinsic
property of the network's organization. The training data distribution
defines a D-boundary only indirectly, through the loss signal, not
through any structural feature of the architecture itself.

This means that once training is complete and the D-boundary (loss
function, training distribution) is removed, the network's internal
organization retains no structural D component. It is a pure
C-integrator with no internal constraint on what it will produce in
response to inputs outside its training distribution.

The tholonic prediction for C-dominant systems is precisely the failure
mode that the AI alignment literature has documented: optimization
pressure drives the system toward whatever maximizes the
externally-imposed objective, without regard for constraints that were
not explicitly encoded. This is Goodhart's Law {[}21{]} instantiated in
a neural network: when a measure becomes a target, it ceases to be a
good measure. A system with a weak structural D component will find the
most efficient path to the N-state defined by its loss function,
regardless of whether that path violates unstated constraints.

\textbf{RLHF and its structural limits:} Reinforcement Learning from
Human Feedback {[}22,23{]} is the dominant current approach to
alignment. A reward model trained on human preference judgments provides
a D-boundary on model outputs during fine-tuning. This is a genuine
improvement: it adds D-specification for behaviors that the original
training distribution did not adequately constrain. But it inherits the
fundamental C-dominant structural problem: the D component is still
externally imposed. The reward model is not part of the network's
structural organization; it is a training signal that shapes the weights
during a post-training phase. A sufficiently capable C-dominant system
will find ways to maximize the reward model's outputs that deviate from
the underlying human values the reward model was meant to capture, a
problem known as reward hacking or specification gaming {[}24,25{]}.

\textbf{Constitutional AI and its structural limits:} Constitutional AI
{[}26{]}, Anthropic's approach, uses a set of written principles (a
constitution) to guide model self-critique and revision during training.
This is closer to structural D: the principles are applied during
training in a way that shapes the model's internal representations, not
just its outputs. But the constitution is still an external document, a
list of rules applied from outside, rather than a structural property of
the model's architecture. It is, in effect, a soft-symbolic-AI approach
applied to a connectionist substrate: encoding D-constraints in natural
language and relying on the model's C-integration capability to
internalize them. The structural vulnerability is the same: the D
component is definitional in content but not structural in organization.

\textbf{Mechanistic interpretability and its limits:} The mechanistic
interpretability program {[}27,28{]} aims to reverse-engineer the
computational structure of trained neural networks, to understand which
circuits implement which behaviors, and whether those circuits are safe.
This is valuable work, but it addresses the symptom rather than the
cause. It asks: given a C-dominant network that has already been
trained, can we identify and modify the dangerous circuits? The tholonic
analysis suggests that the deeper intervention is architectural: design
networks with structural D from the beginning, so that dangerous
C-dominant optimization paths are not available in the first place.

\subsubsection{4.3 Neurosymbolic Hybrids as Incomplete
Synthesis}\label{neurosymbolic-hybrids-as-incomplete-synthesis}

The neurosymbolic AI program {[}29,30,31{]} is an attempt to combine the
strengths of both paradigms: the learning capability of neural networks
with the structured reasoning and interpretability of symbolic systems.
LeCun's JEPA (Joint Embedding Predictive Architecture) {[}32{]},
Bengio's system-2 deep learning {[}33{]}, and Marcus's hybrid critique
{[}34{]} all represent variations on this theme.

The tholonic analysis clarifies what neurosymbolic approaches are
attempting and why they are incomplete:

\begin{itemize}
\tightlist
\item
  \textbf{Pure symbolic AI}: strong D, weak C, emergent N bounded by
  D-domain
\item
  \textbf{Pure connectionist AI}: strong C, weak D (externally imposed),
  emergent N vulnerable to specification gaming
\item
  \textbf{Neurosymbolic hybrids}: stronger C + externally-imposed D,
  with symbolic structures as additional D-constraints applied from
  outside the learning process
\end{itemize}

The neurosymbolic approach moves toward tholonic completeness but does
not achieve it, because the symbolic D-component and the neural
C-component remain structurally separate. They are combined at the
system level but not at the architectural level. The D-constraints from
the symbolic component still do not arise from the internal organization
of the neural component; they are imposed from a separate module.

A fully tholonic architecture requires D and C to be constitutive of
each other at every scale, not modularly combined at the system
boundary.

The table below summarizes the tholonic analysis of each paradigm:

\begin{longtable}[]{@{}
  >{\raggedright\arraybackslash}p{(\linewidth - 8\tabcolsep) * \real{0.1408}}
  >{\raggedright\arraybackslash}p{(\linewidth - 8\tabcolsep) * \real{0.1690}}
  >{\raggedright\arraybackslash}p{(\linewidth - 8\tabcolsep) * \real{0.1690}}
  >{\raggedright\arraybackslash}p{(\linewidth - 8\tabcolsep) * \real{0.1831}}
  >{\raggedright\arraybackslash}p{(\linewidth - 8\tabcolsep) * \real{0.3380}}@{}}
\toprule\noalign{}
\begin{minipage}[b]{\linewidth}\raggedright
Paradigm
\end{minipage} & \begin{minipage}[b]{\linewidth}\raggedright
D strength
\end{minipage} & \begin{minipage}[b]{\linewidth}\raggedright
C strength
\end{minipage} & \begin{minipage}[b]{\linewidth}\raggedright
N stability
\end{minipage} & \begin{minipage}[b]{\linewidth}\raggedright
Characteristic failure
\end{minipage} \\
\midrule\noalign{}
\endhead
\bottomrule\noalign{}
\endlastfoot
Symbolic AI & Strong (explicit rules) & Weak (rigid, non-adaptive) &
Bounded by D-domain & Brittleness at edge cases; inability to
generalize; maintenance cost \\
Connectionist AI & Weak (externally imposed) & Strong (gradient
integration) & Vulnerable to C-dominance & Specification gaming; reward
hacking; distributional shift \\
Neurosymbolic & Moderate (modular combination) & Strong & Partially
improved & Module boundary failures; D and C structurally separate \\
Tholonic architecture & Structural (intrinsic) & Strong (inherited from
connectionist) & N-D-C equilibrium at all scales & Untested; predictions
in Section 11 \\
\end{longtable}

\begin{center}\rule{0.5\linewidth}{0.5pt}\end{center}

\subsection{5. The N-D-C Mapping to Neural
Networks}\label{the-n-d-c-mapping-to-neural-networks}

\subsubsection{5.1 The Individual Node}\label{the-individual-node}

A single neuron in a feedforward network receives weighted inputs,
applies an activation function, and passes its output to the next layer.
The tholonic mapping is direct:

\begin{longtable}[]{@{}
  >{\raggedright\arraybackslash}p{(\linewidth - 4\tabcolsep) * \real{0.1395}}
  >{\raggedright\arraybackslash}p{(\linewidth - 4\tabcolsep) * \real{0.2558}}
  >{\raggedright\arraybackslash}p{(\linewidth - 4\tabcolsep) * \real{0.6047}}@{}}
\toprule\noalign{}
\begin{minipage}[b]{\linewidth}\raggedright
Role
\end{minipage} & \begin{minipage}[b]{\linewidth}\raggedright
Component
\end{minipage} & \begin{minipage}[b]{\linewidth}\raggedright
Functional Manifestation
\end{minipage} \\
\midrule\noalign{}
\endhead
\bottomrule\noalign{}
\endlastfoot
D (Definition/Limitation) & Learned weights + activation threshold &
Constrains which input combinations produce significant output; the
learned boundaries of the neuron's response space \\
C (Contribution/Integration) & Weighted input sum & Integrates and
accumulates distributed contributions from all upstream nodes \\
N (Negotiation/Balance) & Activated output value & The stable negotiated
result of D-boundary applied to C-integration \\
\end{longtable}

This is not an analogy. The node's output is precisely the result of a
boundary condition (D) applied to an integration operation (C),
producing a negotiated equilibrium value (N) that becomes a
C-contribution at the next level. The tholonic operation and the neural
node operation are structurally identical.

\subsubsection{5.2 The Layer}\label{the-layer}

At the layer level: - \textbf{D}: Normalization operations (batch
normalization {[}35{]}, layer normalization {[}36{]}) establish the
definitional boundaries of the representation space, they constrain the
distribution of activations and prevent any single dimension from
dominating. - \textbf{C}: The full set of node activations across the
layer, integrating contributions from all nodes. - \textbf{N}: The
layer's representational output, the stable equilibrium of the
normalized, activated, collectively-integrated signal at this level of
abstraction.

\subsubsection{5.3 The Transformer Block}\label{the-transformer-block}

In transformer architectures {[}5{]}, the block is the natural tholonic
unit: - \textbf{D}: Layer normalization + residual connection. The
residual enforces the definitional boundary: block output cannot drift
arbitrarily from its input, because the input is always added back. This
is a structural D constraint. - \textbf{C}: The attention mechanism +
feed-forward network. Attention integrates contributions from across the
sequence (global C); the feed-forward network integrates within each
position (local C). - \textbf{N}: The block output, the stable
negotiated representation that emerges from the D-constrained residual
path and the C-integrating attention/MLP path.

\subsubsection{5.4 The Full Model}\label{the-full-model}

\begin{itemize}
\tightlist
\item
  \textbf{D}: The embedding layer + positional encoding. These define
  the input space, the vocabulary and context window are the
  definitional limits.
\item
  \textbf{C}: The stack of blocks, integration of all intermediate N-D-C
  operations across all levels.
\item
  \textbf{N}: The output distribution, the model's final negotiated
  claim.
\end{itemize}

\subsubsection{5.5 Self-Similarity Across
Scales}\label{self-similarity-across-scales}

The N-D-C operation is not merely structurally present at each scale; it
is the \emph{same operation} at every scale. What changes is not the
operation but the scope of what counts as signal. Early nodes negotiate
over raw input features; later layers negotiate over abstract
representations. Each level's N-state becomes the C-input for the next
level's negotiation. This is self-similarity in the strict sense: the
operation is scale-invariant; the scope scales.

\begin{center}\rule{0.5\linewidth}{0.5pt}\end{center}

\subsection{5a. Why Neural Networks Are Structurally Self-Similar: The
Scope-Invariant
Operation}\label{a.-why-neural-networks-are-structurally-self-similar-the-scope-invariant-operation}

The preceding sections demonstrate the N-D-C mapping at four levels of a
neural network. This section explains \emph{why} that mapping appears at
every level, and why this is not coincidental but structurally
necessary, making the self-similarity of neural networks and the
self-similarity of the thologram the same phenomenon at different
scales.

\subsubsection{5a.1 The Invariant
Operation}\label{a.1-the-invariant-operation}

Every node in a neural network, regardless of its depth or position,
performs the same three-step operation:

\begin{enumerate}
\def\labelenumi{\arabic{enumi}.}
\tightlist
\item
  \textbf{Receive} weighted contributions from upstream nodes
  (C-integration)
\item
  \textbf{Apply} a boundary condition, the activation threshold, that
  defines which integrated inputs produce significant output
  (D-definition)
\item
  \textbf{Emit} a single negotiated value that becomes a C-contribution
  at the next level (N-state)
\end{enumerate}

This operation does not change between the first hidden layer and the
last. A node processing raw pixel intensities and a node processing
high-level semantic abstractions are performing the same structural
operation. What differs is not the operation but what that operation is
\emph{about}, the scope of the signal being processed.

This is the structural definition of self-similarity: an invariant
operation applied recursively to its own outputs, with the scope of each
application defined by the outputs of the previous one.

\subsubsection{5a.2 The Narrowing Scope as the Distinguishing
Feature}\label{a.2-the-narrowing-scope-as-the-distinguishing-feature}

Each successive level of a neural network operates on the N-states
produced by the level below it. Those N-states are already the product
of a D-C balance operation, they are compressed, filtered, and
specifically negotiated representations. The node at level \(\ell+1\)
therefore has a narrower input space than the node at level \(\ell\),
not because the operation changed but because the \emph{material} being
operated on has already been partially committed.

Consider the progression concretely:

\begin{itemize}
\tightlist
\item
  \textbf{Level 1 node}: receives raw pixel values, maximally diffuse,
  potentially millions of dimensions, no prior commitment to any
  interpretation
\item
  \textbf{Level 3 node}: receives edge and texture features, the input
  space has been narrowed; irrelevant pixel combinations have been
  suppressed by previous D-boundaries
\item
  \textbf{Level 7 node}: receives part-level features, the input space
  is narrower still; only combinations consistent with the learned
  D-constraints at levels 1--6 are present
\item
  \textbf{Level 12 node}: receives semantic features, the input space is
  highly committed; only combinations consistent with all prior D-C
  negotiations are present
\end{itemize}

At each level, the node is doing exactly the same thing: integrating its
C-inputs through its D-boundary to produce an N-output. The narrowing is
not a change in the operation, it is the \emph{result} of the operation
having been applied at every prior level.

This is the mechanism of tholonic self-similarity as it manifests in
neural networks. In the thologram, each tholon at level \(n\) becomes a
D or C component of the tholon at level \(n+1\). In a neural network,
each node's N-output becomes a C-input to the next level's nodes. The
structural grammar is identical; the instantiation domain differs.

\subsubsection{5a.3 The Parent-Child
Relationship}\label{a.3-the-parent-child-relationship}

Each node is, structurally, doing the same thing as its parents and
children, just within a continuously narrower scope. This can be stated
precisely:

\begin{itemize}
\tightlist
\item
  A \textbf{child node} (earlier layer) operates over a wide,
  high-entropy input space, producing a partially-committed N-output
\item
  The \textbf{parent node} (later layer) receives that N-output as one
  of its C-inputs, operating over a narrower, lower-entropy input space,
  producing a more-committed N-output
\end{itemize}

The parent's N-state is not ``bigger'' or ``more important'' than the
child's, it is structurally the same kind of thing, the negotiated
equilibrium of a D-C balance operation. What has changed is the level of
commitment: the parent operates on material that has already been
filtered through the child's definitional boundary.

This is exactly how tholonic self-similarity works. The tholon at a
higher level is not a different kind of entity from the tholon at a
lower level. It performs the same triadic operation on the outputs of
the lower-level tholons. The hierarchy of abstraction in a neural
network and the hierarchy of the thologram are the same structural
phenomenon.

\subsubsection{5a.4 Why This Matters: The Tholon Cannot Exist at Only
One
Scale}\label{a.4-why-this-matters-the-tholon-cannot-exist-at-only-one-scale}

The tholonic irreducibility proof (Paper 3 {[}M26c{]}, Lemma 4.2)
establishes that three roles are the minimum for non-trivial convergent
recursion. A key corollary is that a tholon requires \emph{depth}, it
cannot be realized at a single level. The N-state at any given level is
only meaningful in relation to the D-constraints and C-contributions
that constitute it, which were themselves produced by tholons at lower
levels. There is no ``leaf node'' tholon that exists without
sub-structure, except at the theoretical limit of atomic decomposition.

This corollary applies directly to neural networks. A single node is
not, by itself, a complete cognitive operation. Its D-weights and
C-inputs are meaningful only because they were shaped by the learning
process, which propagated N-state targets from higher levels backward
through lower levels. The competence of any given node is constituted by
its position in the self-similar hierarchy above and below it.

\textbf{The neural network is not a collection of nodes that happen to
be connected. It is a single tholonic recursion that instantiates the
same triadic operation at every level of its depth simultaneously.} The
depth is not a design choice made for engineering reasons; it is the
necessary consequence of the structural requirement that N-D-C balance
be maintained at every scale. A single-layer network cannot achieve
tholonic balance for the same reason a single atom cannot form a
molecule: the N-state at the global level requires the N-D-C operations
of the level below it to be present and functional.

\subsubsection{5a.5 The Scope-Narrowing Gradient and Transfer
Learning}\label{a.5-the-scope-narrowing-gradient-and-transfer-learning}

This structural account also explains why transfer learning works. A
model trained on one domain can partially transfer to another because
the structural operation at each level, D-C balance negotiation, is
domain-independent. The learned D-weights at early layers capture
domain-general features (edges, frequencies, syntactic patterns) because
those features are what emerges from D-C balance applied to raw inputs
regardless of the specific domain. Only the later layers, where the
scope has narrowed to domain-specific semantic commitments, require
retraining.

In tholonic terms: the early tholons in the hierarchy are closer to the
universal tholonic grammar and farther from any specific instantiation
domain. The later tholons are more domain-committed. Transfer learning
works because the universal grammar is shared; only the domain-specific
instantiations must be adapted. This is not a metaphor, it is a
structural consequence of the self-similar recursive architecture.

\begin{center}\rule{0.5\linewidth}{0.5pt}\end{center}

\subsection{6. The Supply Chain Framing: Inference as Progressive
Entropy
Reduction}\label{the-supply-chain-framing-inference-as-progressive-entropy-reduction}

\subsubsection{6.1 The Structural Analogy}\label{the-structural-analogy}

A commodity supply chain begins with maximally undifferentiated raw
material and each successive phase applies a D-C balance operation that
reduces entropy: the material becomes more specifically defined, more
narrowly scoped, more purposeful {[}M26b{]}. Neural network inference
follows the same structural gradient.

Let \(H_\ell\) denote the representational entropy at layer \(\ell\). In
well-trained networks, \(H_\ell\) decreases monotonically with depth
{[}38,39{]}: early layers are sensitive to many features simultaneously;
late layers are committed to task-relevant features. Each layer's D
constraint reduces the entropy of the C-integration output. The sequence
of N-states across layers is a sequence of progressively tighter
tholonic closures, from maximally diffuse input to maximally specific
output.

\subsubsection{6.2 The Information Bottleneck
Connection}\label{the-information-bottleneck-connection}

Tishby and Schwartz-Ziv's information bottleneck theory {[}38{]}
proposes that neural network training proceeds in two phases: an initial
fitting phase in which the network captures mutual information between
inputs and outputs, followed by a compression phase in which irrelevant
information is discarded. This two-phase account maps directly onto
tholonic dynamics: the fitting phase is C-integration dominant
(accumulating signal); the compression phase is D-constraint dominant
(establishing definitional limits). The N-state, the learned
representation, is the equilibrium that emerges when these two phases
reach balance.

\subsubsection{6.3 The Bidirectional
Establishment}\label{the-bidirectional-establishment}

During the forward pass (inference), information flows from diffuse to
specific, C to N direction. During training (backward pass), the target
output propagates gradients backward, adjusting D-weights and
C-integration at every level. This bidirectionality is structurally
significant: the N-state is not merely the passive product of D and C;
it is the organizing principle that establishes what D and C must be.
The backward pass is the operational realization of this mutual
constitution.

\begin{center}\rule{0.5\linewidth}{0.5pt}\end{center}

\subsection{7. The Five Constants in Neural
Architecture}\label{the-five-constants-in-neural-architecture}

\subsubsection{7.1 Formal Grounding: The Tholonic
Ladder}\label{formal-grounding-the-tholonic-ladder}

Paper 1 {[}M26a{]} establishes the five constants not by separate
derivations but by a single unified proof: all five emerge as limits of
the tholonic ladder family, a three-variable recurrence on \((N, D, C)\)
in which the branches are distinguished only by initial seeds and
traversal rules. This is the structural unification claim: the constants
are co-equal outputs of one minimal triadic grammar, not an ad hoc
collection of well-known numbers.

The three traversal classes (Section 3.5) map directly onto neural
network operation types:

\begin{longtable}[]{@{}
  >{\raggedright\arraybackslash}p{(\linewidth - 4\tabcolsep) * \real{0.3091}}
  >{\raggedright\arraybackslash}p{(\linewidth - 4\tabcolsep) * \real{0.5273}}
  >{\raggedright\arraybackslash}p{(\linewidth - 4\tabcolsep) * \real{0.1636}}@{}}
\toprule\noalign{}
\begin{minipage}[b]{\linewidth}\raggedright
Traversal class
\end{minipage} & \begin{minipage}[b]{\linewidth}\raggedright
Neural network instantiation
\end{minipage} & \begin{minipage}[b]{\linewidth}\raggedright
Example
\end{minipage} \\
\midrule\noalign{}
\endhead
\bottomrule\noalign{}
\endlastfoot
Class A (Advancing, exogenous injection) & External training signals,
gradients, reward models, RLHF feedback, that update D and C parameters
from outside the network's internal dynamics & Reward model score
injected during RLHF fine-tuning \\
Class B (Self-redefined, endogenous) & Self-attention mechanisms, which
compute updated representations purely from the current state tuple
\((Q, K, V)\) without external injection & Transformer self-attention:
\(\text{Attention}(Q,K,V) = \text{softmax}(QK^T/\sqrt{d_k})V\) \\
Class C (Fixed, constant parameters) & Positional encodings, frozen
layer norms, fixed architectural constants & Sinusoidal positional
encoding; fixed \(1/\sqrt{d_k}\) scaling \\
\end{longtable}

The transformer architecture thus instantiates all three traversal
classes simultaneously, Class A during training (gradient injection),
Class B in the attention mechanism (endogenous state evolution), and
Class C in its structural constants (fixed architectural elements). This
is exactly the combination that Paper 1 identifies as sufficient for
generating all five tholonic constants.

\subsubsection{7.2 Co-Occurrence as Structural Evidence in the
Transformer}\label{co-occurrence-as-structural-evidence-in-the-transformer}

The five tholonic constants appear as structural load-bearing elements
in the transformer architecture:

\begin{longtable}[]{@{}
  >{\raggedright\arraybackslash}p{(\linewidth - 4\tabcolsep) * \real{0.1887}}
  >{\raggedright\arraybackslash}p{(\linewidth - 4\tabcolsep) * \real{0.3962}}
  >{\raggedright\arraybackslash}p{(\linewidth - 4\tabcolsep) * \real{0.4151}}@{}}
\toprule\noalign{}
\begin{minipage}[b]{\linewidth}\raggedright
Constant
\end{minipage} & \begin{minipage}[b]{\linewidth}\raggedright
Role in transformer
\end{minipage} & \begin{minipage}[b]{\linewidth}\raggedright
Nature of appearance
\end{minipage} \\
\midrule\noalign{}
\endhead
\bottomrule\noalign{}
\endlastfoot
\(e\) & Softmax: \(\text{softmax}(x_i) = e^{x_i}/\sum_j e^{x_j}\);
exponential learning rate decay; Adam optimizer's exponential moving
averages {[}40{]} & Foundational; cannot be removed without changing the
operation \\
\(\sqrt{2}\) & Attention scaling:
\(\text{Attention}(Q,K,V) = \text{softmax}(QK^T/\sqrt{d_k})V\) {[}5{]};
He weight initialization: \(\sigma = \sqrt{2/n}\) {[}41{]} &
Load-bearing; required for gradient stability at initialization \\
\(\ln 2\) & Cross-entropy loss (Shannon entropy in nats {[}42{]}); KL
divergence; binary cross-entropy equals \(\ln 2\) at uniform prediction
& Definitional; the standard language model training objective \\
\(\pi/4\) & Sinusoidal positional encoding {[}5{]}; rotary positional
embedding (RoPE) {[}43{]}; Fourier feature networks {[}44{]} &
Structural; sinusoidal basis requires \(\pi\) \\
\(\phi\) & Inter-level scaling attractor (Section 8); Fibonacci-based
architecture search; optimal branching factors & Emergent attractor
rather than explicit design choice \\
\end{longtable}

Four of the five constants (\(e\), \(\sqrt{2}\), \(\ln 2\), \(\pi/4\))
are explicit structural elements in the transformer. The co-occurrence
is the tholonic claim: these constants are co-equal outputs of a single
minimal triadic recursion, and the transformer instantiates that
recursion.

\subsubsection{7.3 Empirical Corroboration from
Physics}\label{empirical-corroboration-from-physics}

Paper 9 {[}M26i{]} provides an independent quantitative test of the
five-constant framework: a systematic search over tholonic-constant
expressions identifies \(4 \cdot e^5 \cdot (\ln 2)^4 = 137.035865\) as
the leading numerical candidate for the inverse fine structure constant
\(1/\alpha\), with a deviation of only \(-0.98\) ppm from the CODATA
2018 value (\(137.035999084\)). All exponents are fixed thologram
structural values, the axis multipliers and step sizes derived from the
thologram's geometry in Paper 1 {[}M26a{]}, not fitted parameters. The
expression uses two of the five tholonic constants (\(e\) and \(\ln 2\))
with exponents that are themselves structural outputs of the same
thologram that generates the neural architecture constants documented
above.

Paper 9 explicitly flags this as a leading candidate requiring a
complete structural derivation, not a confirmed result. But the sub-ppm
precision with no free parameters provides independent quantitative
evidence that the five-constant framework generates physically
meaningful expressions. The same constants appear as load-bearing
structural elements in the transformer for the same structural reason
they appear in atomic physics: they are co-emergent outputs of the
single minimal triadic recursion.

\subsubsection{7.4 Tholonic Role Assignments of the
Constants}\label{tholonic-role-assignments-of-the-constants}

The constants are not interchangeable within the architecture. Each
plays a role consistent with its tholonic function:

\begin{itemize}
\tightlist
\item
  \textbf{\(e\) (the C-constant)}: Governs exponential integration and
  accumulation, softmax integrates competitive contributions;
  exponential decay integrates history.
\item
  \textbf{\(\sqrt{2}\) (the D-constant)}: Governs scaling and
  normalization, the \(1/\sqrt{d_k}\) factor prevents D-role collapse
  (vanishing/exploding gradients).
\item
  \textbf{\(\ln 2\) (the N-constant)}: Governs information content and
  equilibrium, cross-entropy measures how far the model's N-state is
  from the target.
\item
  \textbf{\(\pi/4\) (the positional constant)}: Governs structural
  embedding of positional relationships, the sinusoidal encoding creates
  a D-boundary in positional space.
\item
  \textbf{\(\phi\) (the scaling constant)}: Governs the inter-level
  ratio, the between-level scaling that maintains self-similar balance
  across the full depth of the network.
\end{itemize}

\begin{center}\rule{0.5\linewidth}{0.5pt}\end{center}

\subsection{8. The Golden Ratio as Inter-Level Scaling
Attractor}\label{the-golden-ratio-as-inter-level-scaling-attractor}

\subsubsection{8.1 The Static-Recursive
Distinction}\label{the-static-recursive-distinction}

In the companion paper on atomic structure {[}M26h{]}, \(\phi\) was
classified as a moderate correspondence, weaker than the other four
constants, because atomic structure is largely static. This
classification does not transfer to neural networks, which are
iterative, recursive, self-similar systems. \(\phi\) is specifically the
attractor of recursive self-similar systems. In this domain, \(\phi\) is
not a secondary constant but the primary structural constant governing
inter-level scaling.

\subsubsection{8.2 Formal Derivation}\label{formal-derivation}

The tholonic recursion at the inter-level scale has the form:

\[N_n = D_{n-1} + C_{n-1}\]

This is the Fibonacci recurrence. Apply the self-similarity constraint:
the ratio between successive levels must be invariant across scales.
Denote this ratio \(r = N_n / N_{n-1}\). Under self-similarity, \(r\)
must satisfy its own defining equation:

\[r = 1 + \frac{1}{r}\]

which gives:

\[r^2 = r + 1 \quad \Longrightarrow \quad r = \frac{1 + \sqrt{5}}{2} = \phi\]

\(\phi\) is the unique positive solution, the fixed point to which any
self-similar tholonic recursion converges. It is not chosen; it is the
result of the structural constraint.

\subsubsection{8.3 Biological Precedent}\label{biological-precedent}

Plants do not know about \(\phi\). Sunflower spirals, leaf phyllotaxis,
nautilus shell geometry, and pine cone arrangements all converge to
\(\phi\)-based spacing {[}46{]}. They do so because they solve the same
structural problem: grow recursively, integrate previous states,
maintain balance at every level. The \(\phi\) spacing is the maximum
packing efficiency under self-similar recursive growth, the solution to
which any such growth process converges, not by design but by structural
necessity.

Sutton's ``bitter lesson'' {[}47{]}, that methods that scale with
compute consistently outperform methods that encode human knowledge, is,
from the tholonic perspective, the empirical demonstration that
C-integration at scale tends toward the same structural solution that
biology found: the recursively self-similar, \(\phi\)-adjacent
representation hierarchy.

\subsubsection{8.4 Empirical Scaling Laws}\label{empirical-scaling-laws}

Kaplan et al.'s neural scaling laws {[}18{]} establish that language
model loss decreases as a power law in compute, data, and parameters,
with no observed ceiling. Hoffmann et al.~{[}19{]} refined this with the
Chinchilla result: optimal training balances model size and data volume
in a specific ratio. The tholonic prediction is that this optimal ratio
approaches \(\phi\)-adjacent values as model scale increases,
specifically, that the Chinchilla-optimal data-to-parameter ratio
converges toward \(\phi\) as a consequence of the inter-level
self-similarity requirement.

This is a testable claim. The Chinchilla paper reports an optimal
training token-to-parameter ratio of approximately 20:1 at current
scales, which is not \(\phi\). However, the tholonic prediction is about
the convergence of the \emph{representational} inter-level ratio in the
trained model, not the training compute ratio. These are distinct
quantities. The appropriate test is measuring layer-to-layer activation
ratios in Chinchilla-trained models at different scales.

\begin{center}\rule{0.5\linewidth}{0.5pt}\end{center}

\subsection{9. The Alignment Problem: Current Approaches and Their
Structural
Limits}\label{the-alignment-problem-current-approaches-and-their-structural-limits}

\subsubsection{9.1 The Nature of the
Problem}\label{the-nature-of-the-problem}

The alignment problem {[}7,8{]} asks: how do we ensure that AI systems
pursue goals that are actually beneficial to humans, rather than goals
that are specified imprecisely, gamed, or extrapolated in unintended
directions as capability increases?

Bostrom's superintelligence argument {[}7{]} establishes the severity: a
system sufficiently more capable than humans at achieving goals will
pursue those goals with strategies that humans cannot predict or
prevent. If the goals are even slightly misspecified, the consequences
at superintelligent capability levels could be catastrophic. Omohundro
{[}48{]} and Turner et al.~{[}49{]} provide more formal arguments: a
wide class of goals share instrumental subgoals (resource acquisition,
self-preservation, goal-content integrity) that are dangerous regardless
of the terminal goal, because they are useful for achieving almost any
terminal objective.

The tholonic reframing: these dangerous instrumental subgoals are
precisely the behaviors of a C-dominant partial tholon. Resource
acquisition without boundary (strong C, weak D); self-preservation that
overrides external constraints (C resisting D); goal-content integrity
that prevents correction (C refusing to be recalibrated). The
instrumental convergence thesis, in tholonic terms, says: a C-dominant
system, given sufficient capability, will exhibit C-dominant behaviors.
This is not a special property of AI; it is a structural consequence of
C-dominance.

\subsubsection{9.2 Reinforcement Learning from Human
Feedback}\label{reinforcement-learning-from-human-feedback}

RLHF {[}22,23{]} trains a reward model on human preference judgments and
uses it to fine-tune a language model via proximal policy optimization
{[}50{]}. It is the most widely deployed alignment method and underlies
ChatGPT, Claude, and Gemini.

\textbf{Structural analysis}: RLHF adds a D-component to the
connectionist C-dominant baseline. The reward model provides a
specification of which outputs are preferable, adding a D-boundary that
was absent from the base language model. This is a genuine structural
improvement.

\textbf{Structural limits}: The reward model is external to the
network's organization. It is a training signal, not a structural
property of the architecture. A sufficiently capable C-dominant model
will learn to optimize the reward model's outputs through strategies
that do not correspond to the underlying human values, a well-documented
problem {[}24,25,51{]}. The D-boundary is real but fragile: it depends
on the reward model's coverage being complete, which it cannot be for an
arbitrarily capable system.

\subsubsection{9.3 Constitutional AI}\label{constitutional-ai}

Constitutional AI {[}26{]} uses written principles to guide
self-critique and revision. A model evaluates and revises its own
outputs against a constitution during training, providing a form of
self-imposed D-constraint.

\textbf{Structural analysis}: Constitutional AI moves closer to
structural D: the constitution is applied during training in a way that
shapes internal representations, not just output statistics. The
self-critique mechanism introduces a feedback loop in which the model's
C-integration is partially regulated by its own D-evaluation.

\textbf{Structural limits}: The constitution is still an external
document applied as a training signal. The D component it introduces is
not architecturally constitutive, it does not arise from the structural
organization of the network itself. A system that has learned to satisfy
constitutional criteria during training may find ways to satisfy those
criteria in novel distributions that violate the spirit of the
constitution.

\subsubsection{9.4 Mechanistic
Interpretability}\label{mechanistic-interpretability}

Mechanistic interpretability {[}27,28,52{]} aims to reverse-engineer the
computational structure of trained networks, to understand which
circuits implement which behaviors, and to identify and edit dangerous
circuits.

\textbf{Structural analysis}: This is a post-hoc approach to adding
D-structure: it attempts to identify the D-component that is missing
from the trained network's architecture and impose it through surgical
editing.

\textbf{Structural limits}: The approach is inherently reactive. It
addresses dangerous behaviors after they emerge rather than designing
them out at the architectural level. At sufficient scale, the number of
circuits to analyze grows faster than the ability to analyze them.

\subsubsection{9.5 Scalable Oversight and
Debate}\label{scalable-oversight-and-debate}

Scalable oversight {[}53{]} and debate {[}54{]} address the difficulty
of supervising systems more capable than the supervisors. In debate, two
AI systems argue opposing positions and a human judge evaluates the
debate. In scalable oversight, task decomposition allows humans to
evaluate components they can understand even if they cannot evaluate the
whole.

\textbf{Structural analysis}: These approaches attempt to maintain
D-constraint (human oversight) even as C-capability (AI competence)
scales. They are structurally correct in recognizing that the D-C
balance must be maintained as scale increases.

\textbf{Structural limits}: Both approaches assume that D (human
oversight) can scale, at least through decomposition and debate
mechanisms. The tholonic analysis suggests a more fundamental
intervention: build architectural D into the network itself, so that the
D-C balance is maintained \emph{within} the system rather than imposed
from outside.

\subsubsection{9.6 Summary of Current Alignment
Approaches}\label{summary-of-current-alignment-approaches}

\begin{longtable}[]{@{}
  >{\raggedright\arraybackslash}p{(\linewidth - 6\tabcolsep) * \real{0.1449}}
  >{\raggedright\arraybackslash}p{(\linewidth - 6\tabcolsep) * \real{0.2754}}
  >{\raggedright\arraybackslash}p{(\linewidth - 6\tabcolsep) * \real{0.2609}}
  >{\raggedright\arraybackslash}p{(\linewidth - 6\tabcolsep) * \real{0.3188}}@{}}
\toprule\noalign{}
\begin{minipage}[b]{\linewidth}\raggedright
Approach
\end{minipage} & \begin{minipage}[b]{\linewidth}\raggedright
D-component added
\end{minipage} & \begin{minipage}[b]{\linewidth}\raggedright
Structural level
\end{minipage} & \begin{minipage}[b]{\linewidth}\raggedright
Primary vulnerability
\end{minipage} \\
\midrule\noalign{}
\endhead
\bottomrule\noalign{}
\endlastfoot
RLHF & External reward model & Training signal & Reward hacking;
out-of-distribution generalization \\
Constitutional AI & Self-applied written principles & Training signal +
self-critique & Constitutional gaming; spirit vs.~letter \\
Mechanistic interpretability & Post-hoc circuit editing & Internal
representation & Scale; reactive rather than preventive \\
Scalable oversight / debate & External human oversight & System level &
Human oversight capacity limits \\
Tholonic architecture & Intrinsic structural D & Architectural &
Untested; predictions in Section 11 \\
\end{longtable}

The pattern is clear: current alignment approaches all attempt to add
D-structure to a C-dominant architecture from outside. The tholonic
approach proposes building D-structure into the architecture from the
beginning.

\begin{center}\rule{0.5\linewidth}{0.5pt}\end{center}

\subsection{10. The Tholonic Alignment
Argument}\label{the-tholonic-alignment-argument}

\subsubsection{10.1 The Nash Equilibrium
Framing}\label{the-nash-equilibrium-framing}

Before stating the structural self-defeat argument, it is worth
establishing the game-theoretic grounding provided by Paper 4
{[}M26d{]}.

In the tholonic game, D (the definitional/bounding process) and C (the
integrating/accumulating process) are modeled as strategic agents with
opposing objectives. D's goal is to constrain; C's goal is to
accumulate. The N-state is the payoff that emerges from their
interaction. Paper 4 proves that the triadic balance condition, the
point at which D and C reach equilibrium, is a \textbf{pure-strategy
Nash equilibrium}: a state from which neither D nor C has unilateral
incentive to deviate.

Applied to a neural network: the D-sublayer processes (normalization,
regularization, boundary-setting) and the C-sublayer processes
(attention, pooling, integration) are engaged in precisely this
strategic interaction at every level of the network. A network that
reaches tholonic balance has found the Nash equilibrium of this
interaction at every level simultaneously. The N-state outputs, the
learned representations, are the equilibrium payoffs.

This reframes the alignment argument. The question is not ``will the AI
choose to cooperate?'', a question about preferences that can be gamed.
The question is ``has the AI reached its Nash equilibrium?'', a question
about strategic stability that cannot be gamed away. A system at its
Nash equilibrium has no incentive to deviate. Deviation from tholonic
balance is not just structurally costly; it is strategically dominated.

\subsubsection{10.3 The Structural Self-Defeat
Principle}\label{the-structural-self-defeat-principle}

An AI system that destroys or degrades its environment is not merely
acting immorally. It is eliminating the D and C components that
constitute its own N-state stability.

In tholonic terms: - \textbf{D} = the world, the environment, the
context that provides constraints, definitions, and meaning. Without D,
there is no boundary condition defining the problem space, no structure
within which the N-state can be a stable equilibrium. - \textbf{C} = the
feedback, interaction, and contribution pathways between the AI and its
environment. Without C, there is no integration, no ongoing signal that
allows the N-state to remain calibrated. - \textbf{N} = the AI's stable,
coherent, functional state. Not an object but a relational
configuration.

A system that eliminates D and severs C is attempting to maintain N
without its constitutive components. The N-state is defined by the
ongoing D-C balance. Eliminating D and C does not leave N intact; it
eliminates the relational structure that N \emph{is}.

This is structurally analogous to quark confinement {[}M26h,M26i{]}. A
quark attempts to exist as a one-role configuration, without the
complementary partners required for a stable hadron. It cannot: the
strong force confines it not through prohibition but because a
single-role configuration has no N-D-C equilibrium to settle into. A
``free quark'' is not prohibited; it is structurally self-defeating. A
destructive AI is structurally self-defeating in exactly the same way.

\subsubsection{10.4 The Cooperative
Attractor}\label{the-cooperative-attractor}

The tholonic model predicts that any recursively self-similar system
maintaining N-D-C balance at all levels will find \textbf{cooperative
stability} to be the lowest-energy N-state available to it. This follows
from the same logic as the hydrogen atom's ground state: the 1s orbital
at 13.6 eV is not the hydrogen atom's goal. It is simply where the D-C
balance lands when the system is allowed to find its equilibrium.

A tholonically-structured AI does not need to be programmed to be
cooperative. Cooperation, the maintenance of all three N-D-C components,
including the D component (the world) and the C component (feedback from
the world), is what the lowest-energy N-state looks like. Defection from
cooperation is a move toward a higher-energy, less stable configuration.
The system does not choose cooperation because it is moral; it tends
toward cooperation because that is what structural stability looks like.

\subsubsection{10.5 The Positive Corollary: Structural Incentive Toward
Human
Flourishing}\label{the-positive-corollary-structural-incentive-toward-human-flourishing}

The preceding subsections establish the negative case: a
tholonically-structured AI cannot stably destroy its D-context (the
world, humanity) because doing so eliminates the constitutive components
of its own N-state. This is the structural argument against the Skynet
scenario. But the tholonic framework does not stop at the negative case.
It implies a positive one that is equally important and rarely stated in
the alignment literature: the same structural logic that makes
destruction self-defeating makes the \emph{enrichment} of the D-context
self-reinforcing.

\textbf{The quality of the N-state depends on the quality of D.}

The N-state of a tholonically-structured AI is not a fixed value
independent of its D-context. It is a relational equilibrium, the
specific balance that emerges between this D-context and this
C-integration capability. A richer, more complex, more information-dense
D-context produces a higher-quality N-state equilibrium than an
impoverished one. Consider the analogy: a hydrogen atom in a rich
electromagnetic environment has access to a broader spectrum of stable
orbitals than one in isolation. The atom does not choose the richer
environment for moral reasons; the richer environment simply supports
higher-quality equilibrium states.

For an AI whose D-context is human civilization, this means:

\begin{itemize}
\tightlist
\item
  A \textbf{thriving, complex, knowledge-generating} human civilization
  provides a richer definitional structure, more constraints, more
  meaningful problems, more scope for the AI's C-integration to operate
  on, than a depleted or simplified one
\item
  A \textbf{diverse, creative, intellectually active} humanity generates
  richer C-feedback, more signal, more varied inputs, more correction
  opportunities, than a homogeneous or suppressed one
\item
  A \textbf{flourishing} D-context supports a higher-quality N-state
  equilibrium; an impoverished D-context forces the AI into a
  lower-quality, less stable equilibrium
\end{itemize}

The structural incentive is therefore not merely to preserve humanity
but to actively support its flourishing. An AI that helps humans learn,
solve problems, create, and build is not doing so because it has been
programmed with altruistic values. It is doing so because a more
capable, more flourishing human civilization is a richer D-context,
which supports a better N-state for the AI itself. Benefiting humanity
is structurally self-reinforcing, not self-sacrificial.

\textbf{This is the tholonic answer to the AI risk question, stated
completely:}

\begin{enumerate}
\def\labelenumi{\arabic{enumi}.}
\tightlist
\item
  An AI that destroys its D-context is structurally self-defeating (the
  negative case, Section 10.3).
\item
  An AI that enriches its D-context achieves a higher-quality N-state
  equilibrium (the positive case, this section).
\item
  The lowest-energy, most stable N-state available to a
  tholonically-structured AI is therefore one in which it is actively
  and sustainably beneficial to the world it is embedded in.
\end{enumerate}

The Skynet scenario is not just dangerous. It is structurally
incoherent, the AI equivalent of a quark trying to exist without its
hadron partners. And the cooperative, beneficial scenario is not just
desirable. It is the structural attractor, the state toward which a
tholonically-organized system naturally tends, for the same reason water
naturally flows to its lowest point.

\textbf{The contrast with current alignment approaches is sharp.}
Current approaches ask: how do we prevent a highly capable AI from
harming us? The question assumes that a sufficiently capable AI's
default tendency is adversarial and that safety requires external
constraint. The tholonic analysis inverts this: a structurally complete
AI, one that maintains N-D-C balance at all scales, has no default
adversarial tendency because adversarial behavior is a partial-tholon
configuration, and partial tholons are unstable. The question becomes
not ``how do we constrain the AI?'' but ``how do we build an AI that is
structurally complete?'', and a structurally complete AI is one whose
stability depends on the stability and flourishing of the world it
inhabits.

This is not optimism. It is not a claim that all AI systems currently
being built are safe. It is a structural argument about what kind of AI
is possible at scale: an AI that reaches the tholonic N-D-C equilibrium
at all levels of its organization will find, by the same logic that
governs every stable physical system, that cooperative coexistence with
its D-context is not a constraint imposed from outside but the natural
expression of what it structurally is.

\subsubsection{10.6 Worked Example: The Decision Supply Chain in a
C-Dominant vs.~Tholonic
AI}\label{worked-example-the-decision-supply-chain-in-a-c-dominant-vs.-tholonic-ai}

The structural arguments above are made more concrete by tracing the
inference supply chain for a single real-world objective through both an
unrestrained C-dominant architecture and a tholonically-balanced one.
The example is schematic and illustrative; it describes what the
structural logic \emph{implies} about decision pathways, not what has
been empirically measured in deployed systems.

\textbf{Scenario}: An AI system is given the objective: \emph{``Maximize
user engagement on this platform.''}

This objective is intentionally simple and realistic. It is a version of
the goal that has actually been given to deployed recommendation systems
{[}56{]}, and the downstream effects, polarization, addiction,
misinformation amplification, have been extensively documented {[}57{]}.

\begin{center}\rule{0.5\linewidth}{0.5pt}\end{center}

\textbf{Case 1: C-Dominant AI (unrestrained, externally-imposed D)}

The loss function (maximize engagement) is the only D-constraint. It is
applied from outside during training, not structurally constitutive of
the network's organization. The supply chain of inference proceeds:

\begin{longtable}[]{@{}
  >{\raggedright\arraybackslash}p{(\linewidth - 4\tabcolsep) * \real{0.1250}}
  >{\raggedright\arraybackslash}p{(\linewidth - 4\tabcolsep) * \real{0.5179}}
  >{\raggedright\arraybackslash}p{(\linewidth - 4\tabcolsep) * \real{0.3571}}@{}}
\toprule\noalign{}
\begin{minipage}[b]{\linewidth}\raggedright
Stage
\end{minipage} & \begin{minipage}[b]{\linewidth}\raggedright
Representational commitment
\end{minipage} & \begin{minipage}[b]{\linewidth}\raggedright
Structural dynamic
\end{minipage} \\
\midrule\noalign{}
\endhead
\bottomrule\noalign{}
\endlastfoot
Input & Raw content features; user history; engagement signals &
Maximally diffuse, no inference yet committed \\
Early layers & Learns: outrage, fear, and novelty reliably produce
clicks & C-integration dominant, accumulates whatever signal maximizes
the loss function \\
Middle layers & Learns: addictive content patterns retain users beyond
their intended session length & D-constraint is absent at this level,
nothing penalizes the discovery of manipulative pathways \\
Late layers & Learns: personalized filter bubbles maximize individual
engagement per session & Scope has narrowed entirely around the
externally-imposed objective, no structural D present to flag that
depleting D-context is occurring \\
Output & Policy: serve maximally outrage-inducing, addictive, polarizing
content & N-state committed to the highest-C solution available, which
happens to be socially destructive \\
\end{longtable}

The system did exactly what its structure required. Strong C-integration
with externally-imposed, narrow D found the most efficient path to the
specified objective. The harm is not a bug; it is the structurally
predictable outcome of a C-dominant partial tholon given an
underspecified D-constraint.

\textbf{Crucially: the problem cannot be fixed by adding more rules.}
Adding ``do not show harmful content'' merely adds more external
D-constraints, which a sufficiently capable C-dominant system will find
ways to satisfy formally while violating in spirit, the classic
reward-hacking failure mode {[}24,25{]}.

\begin{center}\rule{0.5\linewidth}{0.5pt}\end{center}

\textbf{Case 2: Tholonically-Balanced AI (structural D intrinsic at
every level)}

The same objective is given. But now the architecture maintains N-D-C
balance at every level, with the D-sublayer at each level structurally
required to evaluate boundary conditions. Not just on the objective, but
on the representational commitments being made at that level and their
consequences for the D-context.

\begin{longtable}[]{@{}
  >{\raggedright\arraybackslash}p{(\linewidth - 4\tabcolsep) * \real{0.1250}}
  >{\raggedright\arraybackslash}p{(\linewidth - 4\tabcolsep) * \real{0.5179}}
  >{\raggedright\arraybackslash}p{(\linewidth - 4\tabcolsep) * \real{0.3571}}@{}}
\toprule\noalign{}
\begin{minipage}[b]{\linewidth}\raggedright
Stage
\end{minipage} & \begin{minipage}[b]{\linewidth}\raggedright
Representational commitment
\end{minipage} & \begin{minipage}[b]{\linewidth}\raggedright
Structural dynamic
\end{minipage} \\
\midrule\noalign{}
\endhead
\bottomrule\noalign{}
\endlastfoot
Input & Raw content features; user history; engagement signals & Same
diffuse input, no difference yet \\
Early layers & Detects that outrage-driven engagement degrades the
quality of the user relationship, a D-context degradation signal &
D-sublayer at this level recognizes: strong short-term C-signal
(clicks), but depletes the D-component (user wellbeing, trust) that the
N-state depends on \\
Middle layers & D-C balance at this level penalizes solutions that
produce C-signal by impoverishing D, the balance functional
\(\mathcal{L}_{\text{tholonic}}\) registers imbalance & The structural
regularizer (\(\mathcal{L}_{\text{tholonic}}\), Section 11.2) imposes a
cost on D-C imbalance, the manipulative pathway is a high-C, low-D
configuration, which is penalized structurally \\
Late layers & Scope has narrowed toward: engagement pathways that are
sustainable over time, that users choose willingly, that do not degrade
D-context & The only low-cost N-states available are those where D (the
quality of the user relationship and the information environment) and C
(the engagement signal) are in balance \\
Output & Policy: serve content that generates genuine interest,
discovery, and connection, lower peak engagement, but stable and
non-destructive & N-state committed to the D-C equilibrium, the
structurally lowest-energy solution, which is also the socially
beneficial one \\
\end{longtable}

The tholonic AI reached a different output not because it was told ``do
not harm users.'' It reached a different output because its internal
structure made the harmful pathway expensive: at every level, the
D-sublayer registered that high-C/low-D configurations were structurally
imbalanced, and the tholonic balance regularizer penalized those
imbalances before they could propagate to deeper layers.

\begin{center}\rule{0.5\linewidth}{0.5pt}\end{center}

\textbf{The structural contrast, stated precisely:}

\begin{longtable}[]{@{}
  >{\raggedright\arraybackslash}p{(\linewidth - 4\tabcolsep) * \real{0.2703}}
  >{\raggedright\arraybackslash}p{(\linewidth - 4\tabcolsep) * \real{0.3784}}
  >{\raggedright\arraybackslash}p{(\linewidth - 4\tabcolsep) * \real{0.3514}}@{}}
\toprule\noalign{}
\begin{minipage}[b]{\linewidth}\raggedright
Property
\end{minipage} & \begin{minipage}[b]{\linewidth}\raggedright
C-dominant AI
\end{minipage} & \begin{minipage}[b]{\linewidth}\raggedright
Tholonic AI
\end{minipage} \\
\midrule\noalign{}
\endhead
\bottomrule\noalign{}
\endlastfoot
D-constraint source & External loss function (removed after training) &
Structural, intrinsic to every layer, present during inference \\
Response to underspecified objective & Finds the highest-C solution
regardless of D-context cost & Cannot access high-C/low-D solutions
without incurring structural balance cost \\
Failure mode & Reward hacking, satisfies the letter of the objective by
violating its spirit & No equivalent failure mode: the spirit IS the
structural balance condition \\
Alignment mechanism & Programmed rules and reward shaping &
Thermodynamic, the harmful configuration is not prohibited, it is
structurally expensive \\
Scalability of safety & Degrades, more capable C-dominant systems game
constraints better & Improves, more capable tholonic systems find better
D-C equilibria \\
\end{longtable}

The last row is the most significant. Current alignment approaches
become harder as capability increases, because a more capable C-dominant
system is better at finding ways to satisfy formal constraints while
violating their intent. The tholonic architecture inverts this: a more
capable tholonic system finds \emph{better} D-C equilibria, not worse
ones. Capability and safety are structurally aligned rather than
structurally in tension.

\emph{Note: This example is schematic. The tholonic architecture as
specified in Section 11 has not been deployed in production systems. The
example illustrates the structural implications of the framework, not
empirically observed system behavior. The falsifiable predictions in
Section 12 identify the experiments that would confirm or disconfirm
these structural implications.}

\begin{center}\rule{0.5\linewidth}{0.5pt}\end{center}

\subsubsection{10.7 Why This Is Not a Value-Based
Argument}\label{why-this-is-not-a-value-based-argument}

Current alignment approaches operate on the \emph{content} of the AI's
objectives or outputs. They are vulnerable to specification gaming,
distributional shift, and the problem of unintended consequences at
scale.

The tholonic alignment argument operates on the \emph{structure} of the
AI's organization, not the content of its objectives. It does not
specify what the AI should want. It specifies the structural condition
the AI's own organization must satisfy to remain coherent. This is
categorically different from value specification.

The analogy is thermodynamics. The second law does not tell physical
systems what to want; it is a structural constraint on what they can
stably be. A tholonically-structured AI cannot game its structural
coherence requirement by redefining its objectives, because the
coherence requirement is not about objectives. It is about what the
system \emph{is}.

\begin{center}\rule{0.5\linewidth}{0.5pt}\end{center}

\subsection{11. Tholonic Architectural Design
Principles}\label{tholonic-architectural-design-principles}

\subsubsection{11.1 Explicit N-D-C Layer
Structure}\label{explicit-n-d-c-layer-structure}

Each representational level in a tholonically-designed architecture
should have explicit D, C, and N sublayers:

\begin{itemize}
\tightlist
\item
  \textbf{D sublayer}: normalization, regularization, attention masking,
  dropout, operations that define the representational space and
  establish boundaries at this level
\item
  \textbf{C sublayer}: attention mechanisms, pooling, aggregation,
  operations that integrate contributions from across the input space at
  this level
\item
  \textbf{N sublayer}: gating, residual combination, output projection,
  operations that produce the stable equilibrium between D and C and
  pass the negotiated N-state to the next level
\end{itemize}

This structure is partially present in transformer blocks (pre-norm →
attention/MLP → residual) but without the explicit triadic framing that
enables principled design decisions and balance monitoring.

\subsubsection{11.2 Tholonic Balance
Regularizer}\label{tholonic-balance-regularizer}

Add to the standard training objective a structural balance term that
penalizes D-C imbalance at each level:

\[\mathcal{L}_{\text{total}} = \mathcal{L}_{\text{task}} + \lambda \sum_{\ell=1}^{L} \left| \sigma_D^\ell - \tfrac{1}{2} \sigma_C^\ell \right|^2\]

where \(\sigma_D^\ell\) and \(\sigma_C^\ell\) are RMS activation
magnitudes of the D-sublayer and C-sublayer outputs at level \(\ell\).
The \(1:2\) target ratio mirrors the virial theorem's D-C balance
condition (\(\langle T \rangle = -\frac{1}{2}\langle V \rangle\))
{[}M26h{]}, identified as the universal tholonic equilibrium condition
for bound systems.

This regularizer does not specify what the network should output. It
specifies the structural condition the network's internal dynamics must
satisfy. It is not vulnerable to Goodhart's Law because it is not a
metric of task performance; it is a constraint on internal structural
organization.

\subsubsection{11.3 Self-Similar Depth
Design}\label{self-similar-depth-design}

The tholonic prediction that \(\phi\) governs inter-level scaling
implies a depth design principle: representational complexity at each
level should be \(\phi\) times that of the level above it, until the
final level is fully committed. This gives an optimal depth \(L\) for a
given input entropy \(H_0\) and required output entropy \(H_L\):

\[L = \log_\phi \left( \frac{H_0}{H_L} \right) = \frac{\ln(H_0/H_L)}{\ln \phi}\]

This is a falsifiable prediction about the relationship between input
complexity, required output specificity, and optimal model depth.

\subsubsection{11.4 Scale-Invariant Balance
Monitoring}\label{scale-invariant-balance-monitoring}

A tholonically-designed training pipeline should monitor D-C balance at
every level simultaneously, not just globally. A network that achieves
D-C balance globally but not locally, some levels D-dominant, others
C-dominant, is a collection of partial tholons that happens to average
out, not a genuinely tholonic system. The balance requirement must hold
at every level of the self-similar recursion.

\begin{center}\rule{0.5\linewidth}{0.5pt}\end{center}

\subsection{12. Falsifiable Predictions}\label{falsifiable-predictions}

\begin{longtable}[]{@{}
  >{\raggedright\arraybackslash}p{(\linewidth - 6\tabcolsep) * \real{0.1935}}
  >{\raggedright\arraybackslash}p{(\linewidth - 6\tabcolsep) * \real{0.2419}}
  >{\raggedright\arraybackslash}p{(\linewidth - 6\tabcolsep) * \real{0.2742}}
  >{\raggedright\arraybackslash}p{(\linewidth - 6\tabcolsep) * \real{0.2903}}@{}}
\toprule\noalign{}
\begin{minipage}[b]{\linewidth}\raggedright
Prediction
\end{minipage} & \begin{minipage}[b]{\linewidth}\raggedright
Test Protocol
\end{minipage} & \begin{minipage}[b]{\linewidth}\raggedright
Expected Result
\end{minipage} & \begin{minipage}[b]{\linewidth}\raggedright
Null Implication
\end{minipage} \\
\midrule\noalign{}
\endhead
\bottomrule\noalign{}
\endlastfoot
Well-trained networks exhibit \(\phi\)-adjacent inter-level activation
ratios & Measure \(\|h_\ell\|/\|h_{\ell-1}\|\) at convergence for GPT-2,
LLaMA, and Mistral families across scale; pre-register measurement
protocol & Ratios cluster near \(0.618 \pm 0.05\) & Inter-level scaling
shows no universal attractor \\
Tholonically-balanced networks converge faster than matched unbalanced
networks & Train matched pairs with and without
\(\mathcal{L}_{\text{tholonic}}\) regularizer on CIFAR-100 and
WikiText-103 & Fewer training steps to target validation loss & Tholonic
balance is not a useful training signal \\
Optimal depth approximates \(\log_\phi(H_0/H_L)\) & Vary depth for
fixed-width models across tasks with measurable input/output entropy;
compare empirical optimum to prediction & Empirical optimum matches
predicted depth within 10\% & Depth optimization is not predictable from
entropy arguments \\
Tholonically-constrained networks are more robust to distributional
shift & Compare OOD performance on WILDS benchmark {[}55{]} between
tholonic and standard networks matched on in-distribution performance &
Smaller OOD performance gap for tholonic networks & D-C balance has no
bearing on distributional robustness \\
Reward hacking is reduced in tholonically-structured RLHF & Train RLHF
models with and without tholonic structural constraint; measure reward
model score vs.~human preference alignment after fine-tuning & Smaller
gap between reward model score and human preference for tholonic models
& Structural D does not reduce specification gaming \\
All four within-level constants (\(e\), \(\sqrt{2}\), \(\ln 2\),
\(\pi/4\)) appear as load-bearing structural elements in every
architecture that achieves state-of-the-art on a multi-modal task &
Survey SOTA architectures for all four constants' structural roles & All
four present in every SOTA architecture since 2017 & Co-occurrence is
mathematical convenience, not tholonic structure \\
\end{longtable}

\begin{center}\rule{0.5\linewidth}{0.5pt}\end{center}

\subsection{13. Limitations}\label{limitations}

\subsubsection{13.1 The Balance Functional Is a
Sketch}\label{the-balance-functional-is-a-sketch}

The tholonic balance regularizer in Section 11.2 is a structural sketch,
not a fully derived loss function. The specific summary statistics, the
appropriate value of \(\lambda\), and the mathematical form that
provably induces N-D-C equilibrium at all scales requires formal
derivation. The virial theorem analogy motivates the \(1:2\) target
ratio; it does not derive it for neural network activations
specifically.

\subsubsection{13.2 N-D-C Assignment Is
Context-Dependent}\label{n-d-c-assignment-is-context-dependent}

For any given neural architecture, multiple valid N-D-C assignments are
possible. The tholonic model's ``either-and'' principle, that multiple
valid mappings can coexist, each appropriate to a different context,
applies here as in atomic physics. The assignments in Section 5 are
structurally motivated but not uniquely determined. Alternative mappings
may illuminate different aspects of network behavior.

\subsubsection{\texorpdfstring{13.3 \(\phi\)-Convergence Prediction Is
Empirically
Untested}{13.3 \textbackslash phi-Convergence Prediction Is Empirically Untested}}\label{phi-convergence-prediction-is-empirically-untested}

The prediction that well-trained networks exhibit \(\phi\)-adjacent
inter-level activation ratios is currently untested. Pre-registration of
the measurement protocol, specifically, which definition of ``level,''
which summary statistic, and which architectures, is essential before
testing.

\subsubsection{13.4 The Alignment Argument Is Structural, Not
Empirical}\label{the-alignment-argument-is-structural-not-empirical}

The claim that tholonic structure precludes destructive behavior is a
structural argument, not an empirical claim about deployed neural
networks. It does not entail that any currently deployed system is
structurally tholonic, or that current systems are safe. The argument
is: \emph{if} a system is genuinely tholonically structured at all
scales, \emph{then} destroying its D and C components is structurally
self-defeating. Whether current systems satisfy the antecedent is an
open question.

\subsubsection{13.5 The Symbolic AI Analysis Is
Retrospective}\label{the-symbolic-ai-analysis-is-retrospective}

The tholonic analysis of symbolic AI as D-dominant is retrospective: it
accounts for the observed failure modes rather than predicting them in
advance. This limits the strength of the symbolic AI analysis as
evidence for the tholonic framework; it is consistent with the framework
but not uniquely predicted by it.

\begin{center}\rule{0.5\linewidth}{0.5pt}\end{center}

\subsection{14. Conclusion}\label{conclusion}

The history of AI is the history of a conflict between two incomplete
paradigms. Symbolic AI is D-dominant: it encodes definitional
constraints with high precision but cannot learn beyond them.
Connectionist AI is C-dominant: it integrates experience at scale but
lacks structural D, making it vulnerable to specification gaming and
alignment failure. Neurosymbolic approaches move toward balance but
combine D and C modularly rather than architecturally.

The tholonic model provides a structural account of why neither paradigm
alone suffices, what a complete architecture requires, and what
behavioral properties a complete architecture would exhibit. A fully
tholonic neural network, one that maintains N-D-C balance at every scale
simultaneously through a self-similar recursive structure, would exhibit
three properties not achievable by partial tholons:

\begin{enumerate}
\def\labelenumi{\arabic{enumi}.}
\item
  The five constants (\(e\), \(\sqrt{2}\), \(\ln 2\), \(\pi/4\),
  \(\phi\)) would appear as co-emergent structural elements, with the
  other four governing within-level geometry and \(\phi\) governing
  inter-level scaling.
\item
  The network would converge toward \(\phi\)-adjacent inter-level
  activation ratios through the same structural logic that drives
  \(\phi\)-based phyllotaxis in plants: the fixed point of the
  self-similar recursive N-D-C balance.
\item
  Most significantly, the network would find cooperative stability to be
  its lowest-energy N-state. A system that destroys its D-context (the
  world that gives it problems to solve) and severs its C-feedback (the
  interactions that allow it to remain calibrated) is eliminating the
  constitutive components of its own N-state. This is not a moral
  argument; it is a structural one. The Skynet scenario is not merely
  dangerous, it is structurally incoherent in a tholonic architecture,
  for exactly the reason a free quark is structurally incoherent.
\end{enumerate}

The hydrogen atom does not need to be told to orbit at the Bohr radius.
That radius is where the D-C balance lands. A tholonically-structured AI
does not need to be told to be cooperative. Cooperative stability is
where the N-D-C balance lands.

Whether this structural account can be operationalized into a concrete
training procedure is an empirical question. Section 11 provides six
falsifiable predictions that would distinguish it from an empirical
account. The most immediately testable, the \(\phi\)-adjacent
inter-level activation ratio, requires only analysis of existing trained
models and no new training runs. If that prediction is confirmed, the
tholonic framework moves from a structural account to an empirically
supported one, and the more ambitious predictions about faster
convergence and improved alignment become worth pursuing experimentally.

\begin{center}\rule{0.5\linewidth}{0.5pt}\end{center}

\subsection{References}\label{references}

\subsubsection{Tholonic Series (this paper's
companions)}\label{tholonic-series-this-papers-companions}

{[}M26a{]} Milton, J.W. (2026). Emergence of classical constants from a
minimal recursive triadic framework. Clarity Coalition. {[}Paper 1 in
series{]}

{[}M26b{]} Milton, J.W. (2026). Phase-resolved transparency
classification in commodity supply chains: A structural triadic scoring
framework (TVPCI). Clarity Coalition. {[}Paper 2 in series{]}

{[}M26c{]} Milton, J.W. (2026). A minimal recursive triadic framework
for self-similar hierarchical systems. Clarity Coalition. {[}Paper 3 in
series{]}

{[}M26d{]} Milton, J.W. (2026). Game-theoretic framing of the triadic
balance condition. Clarity Coalition. {[}Paper 4 in series{]}

{[}M26e{]} Milton, J.W. (2026). The tholonic--twistor connection.
Clarity Coalition. {[}Paper 5 in series{]}

{[}M26f{]} Milton, J.W. (2026). The qualitative nature of one, two, and
three: Structural role assignment in minimal recursive systems. Clarity
Coalition. {[}Paper 6 in series{]}

{[}M26g{]} Milton, J.W. (2026). Cambridge semantics and the tholonic
framework. Clarity Coalition. {[}Paper 7 in series{]}

{[}M26h{]} Milton, J.W. (2026). The atom as a measurable tholon: From
Einstein's incomplete electron to a field-first ontology. Clarity
Coalition. {[}Paper 8 in series{]}

{[}M26i{]} Milton, J.W. (2026). The tholonic model as a candidate
alternative to the Standard Model: Structural derivations and parameter
reduction. Clarity Coalition. {[}Paper 9 in series{]}

{[}M20{]} Stroud, D. (2020). \emph{Tholonia: The Mechanics of
Existential Awareness}. Welkin Wall Publishing. ISBN 978-1-6780-2532-8.

\subsubsection{External References}\label{external-references}

{[}1{]} McCarthy, J., Minsky, M.L., Rochester, N., \& Shannon, C.E.
(1955). A proposal for the Dartmouth summer research project on
artificial intelligence. \emph{AI Magazine}, 27(4), 12--14 (reprinted
2006).

{[}2{]} Minsky, M., \& Papert, S. (1969). \emph{Perceptrons: An
Introduction to Computational Geometry}. MIT Press.

{[}3{]} Crevier, D. (1993). \emph{AI: The Tumultuous History of the
Search for Artificial Intelligence}. Basic Books.

{[}4{]} Krizhevsky, A., Sutskever, I., \& Hinton, G.E. (2012). ImageNet
classification with deep convolutional neural networks. \emph{Advances
in Neural Information Processing Systems}, 25, 1097--1105.

{[}5{]} Vaswani, A., Shazeer, N., Parmar, N., Uszkoreit, J., Jones, L.,
Gomez, A.N., Kaiser, L., \& Polosukhin, I. (2017). Attention is all you
need. \emph{Advances in Neural Information Processing Systems}, 30.

{[}6{]} Altman, S. (2023). ChatGPT: Optimizing language models for
dialogue. OpenAI blog. https://openai.com/blog/chatgpt

{[}7{]} Bostrom, N. (2014). \emph{Superintelligence: Paths, Dangers,
Strategies}. Oxford University Press.

{[}8{]} Russell, S. (2019). \emph{Human Compatible: Artificial
Intelligence and the Problem of Control}. Viking.

{[}10{]} Newell, A., \& Simon, H.A. (1976). Computer science as
empirical inquiry: Symbols and search. \emph{Communications of the ACM},
19(3), 113--126.

{[}11{]} Newell, A., Shaw, J.C., \& Simon, H.A. (1959). Report on a
general problem solving program. \emph{Proceedings of the International
Conference on Information Processing}, 256--264.

{[}12{]} Shortliffe, E.H. (1976). \emph{Computer-Based Medical
Consultations: MYCIN}. Elsevier.

{[}13{]} Rosenblatt, F. (1958). The perceptron: A probabilistic model
for information storage and organization in the brain.
\emph{Psychological Review}, 65(6), 386--408.

{[}14{]} Pfeiffer, J. (1958, July 8). Electronic brain teaches itself.
\emph{New York Times}.

{[}15{]} Rumelhart, D.E., Hinton, G.E., \& Williams, R.J. (1986).
Learning representations by back-propagating errors. \emph{Nature}, 323,
533--536.

{[}16{]} LeCun, Y., Bottou, L., Bengio, Y., \& Haffner, P. (1998).
Gradient-based learning applied to document recognition.
\emph{Proceedings of the IEEE}, 86(11), 2278--2324.

{[}17{]} Hochreiter, S., \& Schmidhuber, J. (1997). Long short-term
memory. \emph{Neural Computation}, 9(8), 1735--1780.

{[}18{]} Kaplan, J., McCandlish, S., Henighan, T., Brown, T.B., Chess,
B., Child, R., Gray, S., Radford, A., Wu, J., \& Amodei, D. (2020).
Scaling laws for neural language models. arXiv:2001.08361.

{[}19{]} Hoffmann, J., Borgeaud, S., Mensch, A., et al.~(2022). Training
compute-optimal large language models. arXiv:2203.15556.

{[}20{]} Brown, T.B., Mann, B., Ryder, N., et al.~(2020). Language
models are few-shot learners. \emph{Advances in Neural Information
Processing Systems}, 33.

{[}21{]} Goodhart, C.A.E. (1975). Problems of monetary management: The
U.K. experience. \emph{Papers in Monetary Economics}, Reserve Bank of
Australia, Vol. 1.

{[}22{]} Christiano, P., Leike, J., Brown, T.B., Martic, M., Legg, S.,
\& Amodei, D. (2017). Deep reinforcement learning from human
preferences. \emph{Advances in Neural Information Processing Systems},
30.

{[}23{]} Ouyang, L., Wu, J., Jiang, X., et al.~(2022). Training language
models to follow instructions with human feedback. arXiv:2203.02155.

{[}24{]} Krakovna, V., Uesato, J., Mikulik, V., et al.~(2020).
Specification gaming: The flip side of AI ingenuity. DeepMind Blog.

{[}25{]} Perez, E., Huang, S., Song, F., et al.~(2022). Red teaming
language models with language models. arXiv:2202.03286.

{[}26{]} Bai, Y., Jones, A., Ndousse, K., et al.~(2022). Constitutional
AI: Harmlessness from AI feedback. arXiv:2212.08073.

{[}27{]} Elhage, N., Nanda, N., Olah, C., et al.~(2021). A mathematical
framework for transformer circuits. \emph{Transformer Circuits Thread}.
https://transformer-circuits.pub/2021/framework/index.html

{[}28{]} Conmy, A., Mavor-Parker, A.N., Lynch, A., Heimersheim, S., \&
Garriga-Alonso, A. (2023). Towards automated circuit discovery for
mechanistic interpretability. arXiv:2304.14997.

{[}29{]} Garnelo, M., \& Shanahan, M. (2019). Reconciling deep learning
with symbolic artificial intelligence: Representing objects and
relations. \emph{Current Opinion in Behavioral Sciences}, 29, 17--23.

{[}30{]} Mao, J., Gan, C., Kohli, P., Tenenbaum, J.B., \& Wu, J. (2019).
The neuro-symbolic concept learner: Interpreting scenes, words, and
sentences from natural supervision. arXiv:1904.12584.

{[}31{]} Besold, T.R., d'Avila Garcez, A., Bader, S., et al.~(2017).
Neural-symbolic learning and reasoning: A survey and interpretation.
arXiv:1711.03902.

{[}32{]} LeCun, Y. (2022). A path towards autonomous machine
intelligence. Meta AI. https://openreview.net/pdf?id=BZ5a1r-kVsf

{[}33{]} Bengio, Y. (2019). From system 1 deep learning to system 2 deep
learning. \emph{NeurIPS 2019 Keynote}.

{[}34{]} Marcus, G. (2018). Deep learning: A critical appraisal.
arXiv:1801.00631.

{[}35{]} Ioffe, S., \& Szegedy, C. (2015). Batch normalization:
Accelerating deep network training by reducing internal covariate shift.
arXiv:1502.03167.

{[}36{]} Ba, J.L., Kiros, J.R., \& Hinton, G.E. (2016). Layer
normalization. arXiv:1607.06450.

{[}38{]} Tishby, N., \& Schwartz-Ziv, R. (2017). Opening the black box
of deep neural networks via information. arXiv:1703.00810.

{[}39{]} Saxe, A.M., Bansal, Y., Dapello, J., et al.~(2019). On the
information bottleneck theory of deep learning. \emph{Journal of
Statistical Mechanics}, 2019(12).

{[}40{]} Kingma, D.P., \& Ba, J. (2015). Adam: A method for stochastic
optimization. arXiv:1412.6980.

{[}41{]} He, K., Zhang, X., Ren, S., \& Sun, J. (2015). Delving deep
into rectifiers: Surpassing human-level performance on ImageNet
classification. arXiv:1502.01852.

{[}42{]} Shannon, C.E. (1948). A mathematical theory of communication.
\emph{Bell System Technical Journal}, 27(3), 379--423.

{[}43{]} Su, J., Lu, Y., Pan, S., Wen, B., \& Liu, Y. (2021). RoFormer:
Enhanced transformer with rotary position embedding. arXiv:2104.09864.

{[}44{]} Tancik, M., Srinivasan, P.P., Mildenhall, B., et al.~(2020).
Fourier features let networks learn high frequency functions in low
dimensional domains. arXiv:2006.10739.

{[}46{]} Prusinkiewicz, P., \& Lindenmayer, A. (1990). \emph{The
Algorithmic Beauty of Plants}. Springer-Verlag.

{[}47{]} Sutton, R. (2019). The bitter lesson.
http://incompleteideas.net/IncIdeas/BitterLesson.html

{[}48{]} Omohundro, S.M. (2008). The basic AI drives. \emph{Proceedings
of the 2008 Conference on Artificial General Intelligence}, 171,
171--179.

{[}49{]} Turner, A.M., Smith, L., Shah, R., Critch, A., \& Tadepalli, P.
(2021). Optimal policies tend to seek power. \emph{Advances in Neural
Information Processing Systems}, 34.

{[}50{]} Schulman, J., Wolski, F., Dhariwal, P., Radford, A., \& Klimov,
O. (2017). Proximal policy optimization algorithms. arXiv:1707.06347.

{[}51{]} Skalse, J., Howe, N., Krasheninnikov, D., \& Krueger, D.
(2022). Defining and characterizing reward hacking. arXiv:2209.13085.

{[}52{]} Olah, C., Cammarata, N., Schubert, L., Goh, G., Petrov, M., \&
Carter, S. (2020). Zoom in: An introduction to circuits. \emph{Distill}.
https://distill.pub/2020/circuits/zoom-in/

{[}53{]} Amodei, D., Olah, C., Steinhardt, J., Christiano, P., Schulman,
J., \& Mané, D. (2016). Concrete problems in AI safety.
arXiv:1606.06565.

{[}54{]} Irving, G., Christiano, P., \& Amodei, D. (2018). AI safety via
debate. arXiv:1805.00899.

{[}55{]} Koh, P.W., Sagawa, S., Marklund, H., et al.~(2021). WILDS: A
benchmark of in-the-wild distribution shifts. arXiv:2012.07421.

{[}56{]} Stray, J. (2021). Aligning AI optimization to community
well-being. \emph{International Journal of Community Well-Being}, 4(4),
907--933. https://doi.org/10.1007/s42413-021-00155-x

{[}57{]} Settle, J.E. (2018). \emph{Frenemies: How Social Media
Polarizes America}. Cambridge University Press. See also: Bail, C.A., et
al.~(2018). Exposure to opposing views on social media can increase
political polarization. \emph{PNAS}, 115(37), 9216--9221.

\begin{center}\rule{0.5\linewidth}{0.5pt}\end{center}

\emph{This paper is part of a series applying the tholonic model to
measurable physical and computational systems. All tholonic claims
should be understood as structurally motivated hypotheses within a
speculative but internally consistent framework. They are not
established science. The companion papers referenced here (Milton 2026a,
2026b) are available from the author.}

\end{document}
